\section{The Fast Causal Inference Algorithm (with input nodes, no selection bias)}\label{sec:fci}

In this final chapter, we will present the Fast Causal Inference (FCI) algorithm \cite{SMR1995,SMR1999}, a highlight in the field of constraint-based causal discovery.
It was originally designed for purely observational (non-experimental) data and relied on the assumption of acyclicity, but it allowed for selection bias due to conditioning on latent variables.
It has a high complexity, but it is of special interest because it can be shown to be \emph{complete} in a certain sense \cite{Zhang2008_AI}. 
While FCI was originally designed for the acyclic setting, it was recently discovered that it also works in case cycles are present (more specifically, for $\sigma$-faithful simple SCMs) \cite{MooijClaassen_UAI_20}.
Furthermore, it has been extended to deal with prior knowledge regarding exogeneity and unconfoundedness of context variables (of the same type that we used for modeling the treatment variable in randomized controlled trials) \cite{Mooij++_JMLR_20}.

We provide here a treatment of FCI that is unique in the sense that it is self-contained, it allows for cycles by assuming that the data was generated by a simple iSCM, and it allows for the presence of exogenous input nodes.
While the extended FCI of \cite{SMR1999} can deal with selection bias, in our treatment here we will assume that no selection bias is present.\footnote{Deriving a version of FCI that allows to deal with cycles, exogenous input nodes and selection bias is an open research problem.}
By allowing for exogenous input nodes, this version of FCI allows one to generalize the idea of causal discovery in a randomized controlled trial to multiple treatment and outcome variables.

\Joris{Think of a better terminology than adding `C' in front of everything. Apparently `CPAG' was already claimed by Richardson for `cyclic' PAG.}

\Joris{Beware! Faithfulness assumption may be different than for JCI! As we are now conditioning on contexts,
it seems we require faithfulness within each conditioned-context. That is, if we have $C1,C2$ we want
faithfulness of $P(X|C1,C2)$, $P(X,C1|C2)$, $P(X,C2|C1)$. Or something like that...}

\Joris{Perhaps we can do all of this without exogenous input nodes by instead looking at $P_{M}(X_O \mid X_S \in \xi_S, X_J)$.
Then we can treat the $J$ nodes as output variables, and still make use of qualitatively similar causal background knowledge
on the relationship between $J$ nodes the other nodes. For example, we might get undirected edges $j_1 \tut j_2$ in the PAG
if $j_1$ and $j_2$ have a common descendant in $S$. We can still try to deduce whatever we can deduce from the conditional
independence model. The difference between boxes and circles then just becomes the difference between $\forall x_J$ and 
$P(X_J)$-a.s. 

I believe that the crux is that if we condition on the $J$ nodes, \emph{and} we know that they can only be non-colliders, then we can ignore their inducing paths. Indeed, if there are any inducing paths between them, they would be blocked anyway. 
So if we can no longer assume them to be non-colliders, there seem to be two options: (i) treat the sub-PAG on $J$ as unknown (i.e., there might be edges, and with arbitrary ancestral relations between the $J$ nodes); and (ii) treat the $J$ nodes as unconnected and pay for that by getting additional inducing paths downstream and by loosing the causal interpretation of the $J$ nodes.

Idea: study RCT in the presence of selection bias.}
\Joris{Reading off confoundedness seems to have changed from FCI-JCI123r to FCI-CCI12!}

\subsection{Inducing walks}

The key notion in this chapter is that of \emph{$\sigma$-inducing walks}.
We define $\sigma$-inducing walk as a generalization to CDMGs of the notion of inducing path \cite{VermaPearl1990} in DAGs. %(primitive inducing path in \cite{RichardsonSpirtes02}).
\begin{Def}\label{def:sigma-inducing_path}
Let $G$ be a CDMG with input vertices $J$, output vertices $V$ and let $i,j \in V \dcup J$ be distinct nodes.
A walk $\pi$ in $G$ between $i$ and $j$ is called \emph{$\sigma$-inducing} if each collider on $\pi$ is in $\Anc_{G}(\{i,j\})$, and each non-endpoint non-collider on $\pi$ is unblockable. 
  If it is a path, it is called an \emph{$\sigma$-inducing path} between $i$ and $j$.\footnote{In most of the literature, the graph is assumed to be acyclic, and then the notion is referred to simply as ``inducing''.}
\end{Def}
If two nodes are adjacent, any edge connecting the two is a $\sigma$-inducing walk between them.
There exists no $\sigma$-inducing walk between two input nodes.
Figure~\ref{fig:sigma-inducing_path} shows some simple nontrivial examples of $\sigma$-inducing paths.

\begin{figure}\centering
  \begin{tikzpicture}
    \begin{scope}
      \node at (-4,0) {(a)};
      \node[ndout] (X1) at (-3,0) {$i$};
      \node[ndout] (X2) at (-1.5,0) {$j$};
      \node[ndout] (X3) at (0,0) {$k$};
      \draw[arout] (X1) edge (X2);
      \draw[arout] (X2) edge (X3);
      \draw[arlat,bend left] (X2) edge (X3);
    \end{scope}
    \begin{scope}[xshift=5.5cm]
      \node at (-4,0) {(b)};
      \node[ndout] (X1) at (-3,0) {$i$};
      \node[ndout] (X2) at (-1.5,0) {$j$};
      \node[ndout] (X3) at (0,0) {$k$};
      \draw[arout] (X1) edge (X2);
      \draw[arout] (X2) edge (X3);
      \draw[arout,bend left] (X3) edge (X2);
    \end{scope}
    \begin{scope}[xshift=11cm]
      \node at (-4,0) {(c)};
      \node[ndout] (X1) at (-3,0) {$i$};
      \node[ndout] (X2) at (-1.5,0) {$j$};
      \node[ndout] (X3) at (0,0) {$k$};
      \node[ndout] (X4) at (-1.5,-1) {$l$};
      \draw[arout] (X1) edge (X2);
      \draw[arout] (X2) edge (X3);
      \draw[arout] (X3) edge (X4);
      \draw[arout] (X4) edge (X1);
    \end{scope}
    \begin{scope}[yshift=-2cm]
      \node at (-4,0) {(d)};
      \node[ndint] (X1) at (-3,0) {$i$};
      \node[ndint] (X2) at (-1.5,0) {$j$};
      \node[ndout] (X3) at (0,0) {$k$};
      \draw[arout] (X2) edge (X3);
    \end{scope}
  \end{tikzpicture}
  \caption{Four examples of non-trivial $\sigma$-inducing paths in CDMGs. (a) The path $i \tuh j \huh k$ is a $\sigma$-inducing path between $i$ and $k$. (b) The path $i \tuh j \tuh k$ is a $\sigma$-inducing path between $i$ and $k$. (c) The path $i \tuh j \tuh k$ is a $\sigma$-inducing path between $i$ and $k$. The nodes $i$ and $k$ cannot be $i\sigma$-separated by any subset not containing $i,k$ (indeed, $i \nisPerp k$ and $i \nisPerp k \mid j$ in all three graphs, and additionally $i \nisPerp k \mid l$ and $i \nisPerp k \mid \{j,l\}$ in the graph in (c)). (d) The path $j \tuh k$ is $\sigma$-inducing, while there is no $\sigma$-inducing path between $i$ and $j$, nor between $i$ and $k$.\label{fig:sigma-inducing_path}}
  %(d) The walk $i \oto j \oto k \oto l$ is a $\sigma$-inducing walk (path) between $i$ and $l$. The nodes $i$ and $l$ cannot be $i\sigma$-separated by any subset not containing $i, l$ (indeed, $i \nisPerp l$, $i \nisPerp l \given j$, $i \nisPerp l \given k$, and $i \nisPerp l \given j,k$).\label{fig:sigma-inducing_walk}}
\end{figure}

The notion of $\sigma$-inducing walk has the following important properties.
\begin{Prp}\label{prp:sigma-inducing_path_properties}
  Let $G$ be a CDMG with input vertices $J$, output vertices $V$ and let $i \in V$ and $j \in V \dcup J$ be distinct nodes. 
  Then the following are equivalent:
  \begin{enumerate}[label=(\roman*)]
    \item There is a $\sigma$-inducing path in $G$ between $i$ and $j$;\label{prp:sigma-inducing_path_properties2:i}
    \item There is a $\sigma$-inducing walk in $G$ between $i$ and $j$;\label{prp:sigma-inducing_path_properties2:ii}
    \item $i \nisPerp_G j \given Z$ for all $Z \subseteq (V \cup J) \sm \{i,j\}$;\label{prp:sigma-inducing_path_properties2:iii}
    \item $i \nisPerp_G j \given Z$ for $Z = (\Anc_G(\{i,j\}) \cup J) \sm \{i,j\}$.\label{prp:sigma-inducing_path_properties2:iv}
  \end{enumerate}
\end{Prp}
\begin{proof}
  The proof is similar to that of Theorem 4.2 in \cite{RichardsonSpirtes02}.

  \ref{prp:sigma-inducing_path_properties2:i} $\implies$ \ref{prp:sigma-inducing_path_properties2:ii} is trivial.

  \ref{prp:sigma-inducing_path_properties2:ii} $\implies$ \ref{prp:sigma-inducing_path_properties2:iii}: Assume the existence of a $\sigma$-inducing walk between $i$ and $j$ in $G$. 
  Let $Z \subseteq (V \cup J) \sm\{i,j\}$. 
  Consider all walks in $G$ between $i$ and $j$ with the property that all colliders on it are in $\Anc_{G}({\{i,j\}) \cup Z}$, and each non-endpoint non-collider on it is not in $Z$ or is unblockable.
  Such walks exist, since the $\sigma$-inducing walk is one. 
  Let $\mu$ be such a walk with a minimal number of colliders.
  We show that all colliders on $\mu$ must be in $\Anc_{G}(Z)$. 
  Suppose on the contrary the existence of a collider $k$ on $\mu$ that is not ancestor of $Z$. 
  It is either ancestor of $i$ or of $j$, by assumption.
  If $j \in J$, it cannot be ancestor of $j$, and hence must be ancestor of $i$.
  Otherwise, we can assume it to be ancestor of $i$ without loss of generality.
  Then there is a directed path $\pi$ from $k$ to $i$ in $G$ that does not pass through any node of $Z$. 
%  Let $k$ be the vertex closest to $j$ on $\mu$ that is also on $\pi$. 
  Then the subwalk of $\mu$ between $k$ and $j$ can be concatenated with the directed path $\pi$
  into a walk between $i$ and $j$ that has the property, but has fewer colliders than $\mu$: a contradiction. 
  Therefore, $\mu$ is $\sigma$-open given $Z$.
  %can be extended to a walk that is $\sigma$-open given $Z$, by replacing each collider $k$ on
  %it by a walk $k \to \dots \to z \ot \dots \ot k$ for some closest descendant $z \in Z$ of $k$ (with possibly $z=k$).
  Hence, $i$ and $j$ are $\sigma$-connected given $Z$.

  \ref{prp:sigma-inducing_path_properties2:iii} $\implies$ \ref{prp:sigma-inducing_path_properties2:iv} is trivial.

  \ref{prp:sigma-inducing_path_properties2:iv} $\implies$ \ref{prp:sigma-inducing_path_properties2:i}:
  Suppose that %$i$ and $j$ are $\sigma'$-connected given any $Z \subseteq V \sm \{i,j\}$.
  %In particular, 
  $i$ and $j$ are $\sigma$-connected given $Z = (\Anc_{G}(\{i,j\}) \cup J) \sm \{i,j\}$. 
  Let $\pi$ be a path between $i$ and $\{j\} \cup J$ that is $\sigma$-open given $Z$. % and contains $i$ and $j$ only once.
  The end nodes of $\pi$ must be $i$ and $j$, because $J \sm \{i,j\} \subseteq Z$.
  We show that $\pi$ must be a $\sigma$-inducing path. 
  First, all colliders on $\pi$ are in $\Anc_G(Z)$, but not in $J$, and hence in $\Anc_{G}(\{i,j\})$.
  Second, let $k$ be any non-endpoint non-collider on $\pi$.
  Then there must be a directed subpath of $\pi$ starting at $k$ that ends either at the first collider on $\pi$ next to $k$ or at an end node of $\pi$, and hence $k$ must be in $Z$. 
  Since $\pi$ is $\sigma$-open given $Z$, $k$ must be unblockable.
  Hence, all non-endpoint non-colliders on $\pi$ must be unblockable.
\end{proof}
In words: there is a $\sigma$-inducing path between two nodes in a CDMG (provided they are not both input nodes) if and only if the two nodes cannot be $i\sigma$-separated by any subset of nodes that does not contain either of the two nodes.
%A similar property holds for $d$-inducing walks and paths, but with $i\sigma$-separation replaced by $id$-separation in points (iii) and (iv).
\begin{Rem}
Two input nodes cannot be $i\sigma$-separated by some subset of other nodes.
Indeed, the trivial path $j$ is $\sigma$-open as long as we don't condition on $j$ itself.
On the other hand, by definition, there is no $\sigma$-inducing path between two input nodes.
This is why Proposition~\ref{prp:sigma-inducing_path_properties} does not consider the case $i,j \in J$.
\end{Rem}

The orientations of the outermost edges on a $\sigma$-inducing path contain important information about ancestral relations.
\begin{Lem}\label{lem:sigma-inducing_path_out-of}
  Let $G$ be a CDMG with input vertices $J$, output vertices $V$ and let $i,j \in J \cup V$ be distinct. 
  If there exists a $\sigma$-inducing path between $i$ and $j$ in $G$,
  %and all $\sigma$-inducing paths in $G$ between $i$ and $j$ are out of $i$, then $i \in \Anc_G(j)$.
  and all $\sigma$-inducing paths in $G$ between $i$ and $j$ are out of $j$, then $j \in \Anc_G(i)$.
\end{Lem}
\begin{proof}
  Let $\mu$ be a $\sigma$-inducing path between $i$ and $j$ in $G$. 
  It must be of the form $i \cdots l \hut j$ (with possibly $l = i$). 
  First we show that $l$ cannot be in $\Sc_G(j)$.
  If $l \in \Sc_G(j)$, then let $\pi$ be a directed path in $G$ from $l$ to $j$ that is entirely contained in $\Sc_G(j)$.
  Let $m$ be the node on $\mu$ closest to $i$ that is also on $\pi$ (possibly $m = l$).
  The subpath of $\pi$ between $j$ and $m$ can be concatenated with the subpath of $\mu$ between $m$ and $i$ into a walk between $j$ and $i$.
  This must be a $\sigma$-inducing path between $i$ and $j$ that is into $j$ by construction: contradiction.
  Hence $l$ cannot be in $\Sc_G(j)$.

  If $\mu$ is a directed path all the way to $i$, then clearly, $j \in \Anc_G(i)$.
  Otherwise, it must contain a collider. 
  Let $k$ be the collider on $\mu$ closest to $j$.
  $k$ must be ancestor of $i$ or $j$.
  In the former case, clearly $j \in \Anc_G(i)$.
  In the latter case, all nodes on the subpath of $\mu$ between $j$ and $k$ must be in $\Sc_G(j)$, a contradiction.
\end{proof}

\begin{Lem}\label{lem:sigma-inducing_path_into_into}
  Let $G$ be a CDMG with input vertices $J$, output vertices $V$ and $i, j \in J \cup V$ distinct.
  If there exists a $\sigma$-inducing path between $i$ and $j$ in $G$ into $j$, and $i \notin \Anc_G(j)$, then there exists a $\sigma$-inducing path between $i$ and $j$ in $G$ that is both into $i$ and into $j$.
\end{Lem}
\begin{proof}
  Let $\mu$ be a $\sigma$-inducing path between $i$ and $j$ in $G$ into $j$.
  This rules out $j \in J$.
  If $\mu$ is into $i$, we are done.
  Therefore, suppose it is of the form $i \tuh \cdots \suh j$. 
  It cannot be a directed path, since $i \notin \Anc_G(j)$.
  Therefore, there must be a collider $k$ on $\mu$ such that $\mu$ is of the form $i \tuh \dots \tuh k \hus \cdots \suh j$
  (with the subpath between $i$ and $k$ directed).
  Then $k \in \Anc_G(i)$, and hence all nodes on $\mu$ between $i$ and $k$ must be in $\Sc_G(i)$.
  Let $\pi$ be a directed path in $G$ from $k$ to $i$ that is entirely contained in $\Sc_G(i)$.
  Let $l$ be the node on $\mu$ closest to $j$ that is also on $\pi$ (possibly $l = k$).
  Then $l \ne j$, because otherwise $j \in \Sc_G(i)$, contradicting $i \notin \Anc_G(j)$.
  The non-trivial subpath of $\pi$ between $i$ and $l$ can be concatenated with the non-trivial subpath of $\mu$ between $l$ and $j$ into a walk between $i$ and $j$.
  This must be a $\sigma$-inducing path between $i$ and $j$ that is both into $i$ and into $j$.
\end{proof}

\subsection{Conditional Directed Partial Ancestral Graphs (CDPAGs)}

It is often convenient when performing causal reasoning to be able to represent a set of CDMGs in a compact way. 
For this purpose, \emph{partial ancestral graphs (PAGs)} have been introduced \cite{SMR1999,Zhang2006}.
Here we will assume no selection bias for simplicity, which allows us to restrict ourselves to discussing 
\emph{directed} PAGs (henceforth abbreviated as DPAGs).
We will extend these objects to allow for input nodes, in which case we will refer to them as \emph{conditional directed} PAGs (CDPAGs).

Conditional directed PAGs have two node types, input nodes and output nodes (just like CDMGs).
They also have multiple edge types. In addition to the three edge types for CDMGs
($\tuh$, $\hut$, $\huh$), there are three new edge types involving a circle: $\huo,\ouo,\ouh$.
Each edge $i \sus j$ therefore has two \emph{edge marks}, one for each node, with each edge mark
either a \emph{tail}, \emph{arrowhead} or \emph{circle}. For example, the directed edge $i \tuh j$ has a 
tail at $i$ and an arrowhead at $j$, while the bi-circle edge $i \ouo j$ has two circle edge marks.\footnote{PAGs have more edge types than DPAGs: they can also have 
undirected or circle-tail edges, i.e., edges of the form $\{\tut,\tuo,\out\}$. This means that for PAGs,
all combinations of edge marks may occur.}
Only 6 out of the 9 possible combinations of edge marks can occur in conditional directed PAGs. We will make use
of the ``$\asterisk$'' symbol to denote any of the three edge marks. So the notation $i \sus j$
includes all 6 possible edge types between $i$ and $j$, whereas $i \hus j$ is shorthand for three
possible edge types.
Edges of the form $i \hut j, i \huo j, i \huh j$ are called \emph{into $i$}, and similarly, edges of the form $i \tuh j, i \ouh j, i \huh j$ are called \emph{into $j$}. Edges of the form $i \tuh j$ and $j \hut i$ are called \emph{out of $i$}.


In order to define CDPAGs, we extend the definitions of (directed) walks, (directed) paths and colliders to cover these new edge types.
\begin{Def}\label{def:more_walks}
  Let $H=\lt J,V,E \rt$ be a mixed graph with input nodes $J$, output nodes $V$ and edges $E$ of the types $\{\tuh,\hut,\huo,\huh,\ouo,\ouh\}$.\footnote{Formally, we no longer introduce separate sets to represent the edges of each type, but merge them into the single set $E$.}
  Let $v,w \in V \cup J$.
  \begin{enumerate}
    \item If there is an edge $v \sus w$ between $v$ and $w$ (of any type), we call $v$ and $w$ \emph{adjacent} in $H$.
    \item Define $\Adj_H(v)$ to be the vertices adjacent to $v$ in $H$.
    \item A triple of distinct nodes $\lt a,b,c \rt$ in $H$ form a \emph{triangle} if each pair of nodes in the triple is adjacent in $H$.
    \item A triple of distinct nodes $\lt a,b,c \rt$ in $H$ is called \emph{unshielded} if $b$ is adjacent to both $a$ and $c$ in $H$, but $a$ is not adjacent to $c$ in $H$.
    \item A \emph{walk between $v$ and $w$} in $H$ is a finite alternating sequence of nodes and edges 
        \[v=v_0, a_0, v_1, \dots v_{n-1}, a_{n-1}, v_n=w\] 
        in $H$ for some $n \ge 0$, i.e.\ such that for every $k=0,\dots,n-1$ 
        we have that $v_{k},v_{k+1} \in V \cup J$ and $a_k = v_k \sus v_{k+1} \in E$, and with end nodes $v_0=v$ and $v_n=w$.
    \item A walk is called a \emph{path} if no node occurs more than once.
    \item A path $v_0 \sus \dots \sus v_n$ in $H$ is called a \emph{possibly directed path from $v_0$ to $v_n$} if for each $i=1,\dots,n$, the edge $v_{i-1} \sus v_i$ is not into $v_{i-1}$ (i.e., the edge is of the form $v_{i-1} \ouo v_i$, $v_{i-1} \ouh v_i$, or $v_{i-1} \tuh v_i$). 
    \item A \emph{directed walk} (path) from $v$ to $w$ in $H$ is a walk (path) of the form:
        \[v=v_0 \tuh v_1 \tuh  \cdots \tuh v_{n-1} \tuh v_n=w,\]
        for some $n \ge 0$.
    \item A \emph{directed cycle} is a directed walk $v \tuh \dots \tuh w$, concatenated with
      the directed edge $w \tuh v$.
    \item An \emph{almost directed cycle} is a directed walk $v \tuh \dots \tuh w$, concatenated
      with the bidirected edge $w \huh v$.
    \item A path $v_0 \sus \dots \sus v_n$ in $H$ is called \emph{uncovered} if every subsequent triple $\lt v_{k-1}, v_k, v_{k+1} \rt$ ($1 < k < n$) is unshielded.
    \item A triple of consecutive nodes $v_{k-1} \sus v_k \sus v_{k+1}$ on a walk in $H$ is called \emph{definite collider}
      if it is of the form $v_{k-1} \suh v_k \hus v_{k+1}$.
    \item A triple of consecutive nodes $v_{k-1} \sus v_k \sus v_{k+1}$ on a walk in $H$ is called a \emph{definite non-collider}
      if it is of the form $v_{k-1} \sus v_k \tuh v_{k+1}$ or $v_{k-1} \hut v_k \sus v_{k+1}$.
    \item If there is a directed walk from $v$ to $w$ in $H$ then we say that \emph{$v$ is ancestor of $w$ in $H$}, 
      and we write $v \in \Anc_H(w)$. For $W \subseteq V \cup J$, we define $\Anc_H(W) := \bigcup_{w\in W} \Anc_H(w)$.
    \item If there is a directed walk from $v$ to $w$ in $H$ then we say that \emph{$w$ is descendant of $v$ in $H$}, 
      and we write $w \in \Desc_H(v)$. For $W \subseteq V \cup J$, we define $\Desc_H(W) := \bigcup_{w\in W} \Desc_H(w)$.
%    \item A non-endpoint non-collider on a walk in $H$ such that all the neighboring nodes on the walk that it points to with a directed edge lie in the same strongly connected component of $H$ is called an \emph{unblockable non-collider}. \Joris{Perhaps only do the $id$-separation stuff, as they will be acyclic?}
%    \item A walk between $v,w \in V \cup J$ (with $v \ne w$) in $H$ is called \emph{$d$-inducing} if every collider on the walk is in $\Anc_H(\{v,w\})$, and it contains no non-endpoint non-collider. \Joris{What to do with circle edge marks?}
%    \item A walk between $v,w \in V \cup J$ (with $v \ne w$) in $H$ is called \emph{$\sigma$-inducing} if every collider on the walk is in $\Anc_H(\{v,w\})$, and every non-endpoint non-collider on the walk is unblockable. \Joris{What to do with circle edge marks?}
    \item A walk between $v,w \in V \cup J$ (with $v \ne w$) in $H$ is called \emph{definitely inducing} if every non-endpoint node is a definite collider in $\Anc_H(\{v,w\})$.
    \item A walk between $v$ and $w$ in $H$ is called \emph{definitely open given $Z \subseteq V \cup J$} if
      \begin{itemize}
        \item every node on the walk is a definite collider or definite non-collider;
        \item every collider is in $\Anc_H(Z)$, and
        \item every non-endpoint non-collider is not in $Z$, and
        \item every endpoint is not in $Z$.
      \end{itemize}
  \end{enumerate}
\end{Def}

We can now define:
\begin{Def}\label{def:dpag}
%  We call a graph $G = \lt V,E\rt$ with nodes $V$ and edges $E$ between ordered node pairs of the types $\{\ttt,\ttc,\to,\ot,\otc,\oto,\ctt,\ctc,\cto\}$ a \emph{partial ancestral graph (PAG)} if
  %It is called a \emph{directed partial ancestral graph (DPAG)} if it only has edges of the types $\{\to,\ot,\oto,\otc,\cto,\ctc\}$.
  A mixed graph $H = \lt J,V,E\rt$ with input nodes $J$, output nodes $V$ and edges $E$ of the types $\{\tuh,\hut,\huo,\huh,\ouo,\ouh\}$ is called a \emph{conditional directed partial ancestral graph (CDPAG)} if all of the following conditions hold:
  \begin{enumerate}
    \item Between any two distinct nodes there is at most one edge, and there are no edges between a node and itself;
    \item The graph contains no directed cycles and no almost directed cycles (``ancestral'');
    \item There is no definitely inducing path between any two non-adjacent nodes (``maximal'');
    \item No input node is adjacent to any other input node;
    \item If there is an edge $j \sus v$ with $j \in J, v \in V$ then it must be of the form $j \tuh v$.
  \end{enumerate}
\end{Def}

CDPAGs are used to represent a set of CDMGs as follows.
\begin{Def}\label{def:dpag_represents_dmg}
  Let $H = \lt J,V,E\rt$ be a mixed graph with input nodes $J$, output nodes $V$ and edges $E$ of the types $\{\tuh,\hut,\huo,\huh,\ouo,\ouh\}$.
  Let $G$ be a CDMG.
  We say that \emph{$H$ represents $G$} if all of the following hold:
  \begin{enumerate}
    \item $G$ and $H$ have the same input nodes $J$ and the same output nodes $V$;
    \item Between any two distinct nodes there is at most one edge;
    \item Two vertices $i,j \in V \cup J$ are adjacent in $H$ if and only if $i \ne j$ and there is a $\sigma$-inducing path between $i$ and $j$ in $G$;
    \item If $i \suh j$ in $H$ (i.e., $i \tuh j$ in $H$ or $i \ouh j$ in $H$ or $i \huh j$ in $H$), then $j \notin \Anc_{G}(i)$;
    \item If $i \tuh j$ in $H$ then $i \in \Anc_{G}(j)$;
    \item If $j \sus v$ in $H$ with $j \in J, v \in V$ then it must be of the form $j \tuh v$.
  \end{enumerate}
\end{Def}
Hence, adjacencies represent $\sigma$-inducing paths, arrowheads represent non-ancestorship, and tails represent ancestorship.
Some examples are given in Figure~\ref{fig:dpags}. 
Note in particular that a directed edge in a CDPAG does not necessarily imply a direct causal relationship (for example, the edge $i \tuh k$ in Figure~\ref{fig:dpags}(c)).
Figure~\ref{fig:nonmaximal_ancestral_graph} provides an example of a mixed graph that satisfies all conditions of a CDPAG except the maximality.

\begin{figure}\centering
  \begin{tikzpicture}
    \begin{scope}
      \node at (-2,0) {(a)};
      \node[ndout] (X1) at (-1,0) {$X_1$};
      \node[ndout] (X2) at (1,0) {$X_2$};
      \node[ndout] (X3) at (0,-1) {$X_3$};
      \node[ndout] (X4) at (0,-2.5) {$X_4$};
      \draw[ouh] (X1) edge (X3);
  %    \draw[biarr,bend right] (X1) edge (X3);
      \draw[ouh] (X2) edge (X3);
  %    \draw[biarr,bend left] (X2) edge (X3);
      \draw[arout] (X3) edge (X4);
    \end{scope}
    \begin{scope}[xshift=6.5cm]
      \node at (-4,0) {(b)};
      \node[ndout] (X1) at (-3,0) {$X_1$};
      \node[ndout] (X2) at (-1.5,0) {$X_2$};
      \node[ndout] (X3) at (0,0) {$X_3$};
      \draw[ouo] (X1) edge (X2);
      \draw[ouo] (X2) edge (X3);
    \end{scope}
    \begin{scope}[xshift=6.5cm,yshift=-1.5cm]
      \node[ndout] (X1) at (-3,0) {$X_1$};
      \node[ndout] (X2) at (-1.5,0) {$X_2$};
      \node[ndout] (X3) at (0,0) {$X_3$};
      \draw[ouh] (X1) edge (X2);
      \draw[arout] (X2) edge (X3);
    \end{scope}
    \begin{scope}[yshift=0cm,xshift=12cm]
      \node at (-4,0) {(c)};
      \node[ndout] (X1) at (-3,0) {$i$};
      \node[ndout] (X2) at (-1.5,0) {$j$};
      \node[ndout] (X3) at (0,0) {$k$};
      \draw[arout] (X1) edge (X2);
      \draw[arout] (X2) edge (X3);
%      \draw[biarr,bend left] (X2) edge (X3);
      \draw[arout,bend right] (X1) edge (X3);
    \end{scope}
    \begin{scope}[yshift=-2cm,xshift=12cm]
      \node at (-4,0) {(d)};
      \node[ndout] (X1) at (-3,0) {$i$};
      \node[ndout] (X2) at (-1.5,0) {$j$};
      \node[ndout] (X3) at (0,0) {$k$};
      \draw[arout] (X1) edge (X2);
      \draw[ouo] (X2) edge (X3);
%      \draw[biarr,bend left] (X2) edge (X3);
      \draw[arout,bend right] (X1) edge (X3);
    \end{scope}
    \begin{scope}[yshift=-4cm,xshift=12cm]
      \node at (-4,0) {(e)};
      \node[ndout] (X1) at (-3,0) {$i$};
      \node[ndout] (X2) at (-1.5,0) {$j$};
      \node[ndout] (X3) at (0,0) {$k$};
      \draw[ouo] (X1) edge (X2);
      \draw[ouo] (X2) edge (X3);
%      \draw[biarr,bend left] (X2) edge (X3);
      \draw[ouo,bend right] (X1) edge (X3);
    \end{scope}
    \begin{scope}[yshift=-4cm,xshift=6.5cm]
      \node at (-4,0) {(f)};
      \node[ndint] (X1) at (-3,0) {$i$};
      \node[ndint] (X2) at (-1.5,0) {$j$};
      \node[ndout] (X3) at (0,0) {$k$};
      \draw[arout] (X2) edge (X3);
    \end{scope}
  \end{tikzpicture}
  \caption{Example DPAGs. (a) DPAG representing the Y-structures in Figure~\ref{fig:Ystruct}. 
  (b) Two different DPAGs that both represent all the LCD DMGs from Figure~\ref{fig:LCD}. 
  (c) DPAG that represents the DMGs in Figure~\ref{fig:sigma-inducing_path}(a).
  (d) DPAG that represents both DMGs in Figure~\ref{fig:sigma-inducing_path}(a-b).
  (e) DPAG that represents the three DMGs in Figure~\ref{fig:sigma-inducing_path}(a-c).
  (f) Unique CDPAG that represents the CDMG in Figure~\ref{fig:sigma-inducing_path}(d).\label{fig:dpags}}
\end{figure}

\begin{figure}\centering
  \begin{tikzpicture}
    \begin{scope}
%      \node at (-2,0) {(d)};
      \node[ndout] (X1) at (-1,0) {$i$};
      \node[ndout] (X2) at (-1,2) {$j$};
      \node[ndout] (X3) at (1,2) {$k$};
      \node[ndout] (X4) at (1,0) {$l$};
      \draw[arlat] (X1) edge (X2);
      \draw[arlat] (X2) edge (X3);
      \draw[arlat] (X3) edge (X4);
      \draw[arout] (X2) edge (X4);
      \draw[arout] (X3) edge (X1);
    \end{scope}
  \end{tikzpicture}
  \caption{This mixed graph is not a valid CDPAG, because it is not maximal: 
  it has a $\sigma$-inducing path $i \huh j \huh  k \huh l$ while $i$ and $l$ are non-adjacent.\label{fig:nonmaximal_ancestral_graph}}
\end{figure}

The following property shows that we can read off ancestral relations from a CDPAG that represents a CDMG:
\begin{Lem}\label{lem:dpag_ancestral_relations}
  Let $H$ be a CDPAG that represents a CDMG $G$. 
%Then:
%  \begin{itemize}
%    \item if $i \tuh j$ in $H$ then $i \in \Anc_{G}(j)$;
%    \item if there is a directed path from $i$ to $j$ in $H$ then $i \in \Anc_{G}(j)$.
%  \end{itemize}
  For $i,j \in V \cup J$: $i \in \Anc_{H}(j)$ implies $i \in \Anc_{G}(j)$.
\end{Lem}
\begin{proof}
%  If $i \tuh j$ in $H$ then, by definition, $i \in \Anc_{G}(\{j\})$ and $j \notin \Anc_{G}(\{i\})$.
%  Combining the two gives that $i \in \Anc_G(j)$.
  If $H$ contains a directed path $i = v_0 \tuh \dots \tuh v_n = j$,
  then each directed edge $v_k \tuh v_{k+1}$ along this directed path implies that $v_k \in \Anc_G(v_{k+1})$. 
  By transitivity, therefore, $i \in \Anc_G(j)$. 
\end{proof}
The following lemma shows that the orientation of edges in a CDPAG that represents a CDMG contains information on the orientation of corresponding $\sigma$-inducing paths in the CDMG.
\begin{Lem}\label{lem:sigma-inducing_path_into}
  Let $H$ be a CDPAG that represents a CDMG $G$ with input nodes $J$ and output nodes $V$.
  Let $i,j \in J \cup V$ be distinct. 
  \begin{enumerate}[label=(\roman*)]
    \item If $i \suh j$ in $H$, then there exists a $\sigma$-inducing path in $G$ between $i$ and $j$ that is into $j$. \label{lem:sigma-inducing_path_into2_into}
%    \item If $i \hus j$ in $H$, then there exists a $\sigma$-inducing path in $G$ between $i$ and $j$ that is into $i$.
    \item If $i \huh j$ in $H$, then there exists a $\sigma$-inducing path in $G$ between $i$ and $j$ that is both into $i$ and into $j$. \label{lem:sigma-inducing_path_into2_intointo}
  \end{enumerate}
\end{Lem}
\begin{proof}
  In both cases, there exists a $\sigma$-inducing path between $i$ and $j$ in $G$ because $i$ and $j$ are adjacent in $H$ and $H$ represents $G$.
  If all $\sigma$-inducing paths between $i$ and $j$ in $G$ were out of $j$, then by Lemma~\ref{lem:sigma-inducing_path_out-of},
  $j \in \Anc_G(i)$, contradicting the orientation $i \suh j$ in $H$. 
  Therefore, there must be a $\sigma$-inducing path between $i$ and $j$ in $G$ that is into $j$.
  This shows~\ref{lem:sigma-inducing_path_into2_into}.
%  If $i \huh j$ in $H$, then with similar reasoning, there must be a $\sigma$-inducing path between $i$ and $j$ in $G$ that is into $j$.
  If $i \huh j$ in $H$, then application of Lemma~\ref{lem:sigma-inducing_path_into_into} gives~\ref{lem:sigma-inducing_path_into2_intointo}.
\end{proof}

We will frequently use that every mixed graph (of a certain type) that represents a CDMG must be a valid CDPAG.
\begin{Prp}\label{prp:dpag_represents_dmg_is_valid}
Let $H = \lt J,V,E\rt$ be a mixed graph with input nodes $J$, output nodes $V$ and edges $E$ of the types $\{\tuh,\hut,\huo,\huh,\ouo,\ouh\}$.
If $H$ represents a CDMG $G$, then $H$ is a CDPAG.
\end{Prp}
\begin{proof}
  ``Ancestral''. Suppose that $H$ contains a directed path $i = v_1 \tuh \dots \tuh v_n = j$.
  Then, each directed edge $v_k \tuh v_{k+1}$ along this directed path implies that $v_k \in \Anc_G(v_{k+1})$. 
  By transitivity, $i \in \Anc_G(j)$. 
  If $H$ contained an edge $j \suh i$, this would imply $i \notin \Anc_G(j)$, a contradiction. 
  Hence such edges cannot occur.
  This means that no directed or almost directed cycles occur in $H$.

  ``Maximal''.
  Suppose there is a definitely inducing path in $H$ between two distinct nodes $u,v \in V \cup J$. 
  Every edge $i \sus j$ on $\mu$ corresponds with a $\sigma$-inducing path in $G$ between $i$ and $j$.
  By Lemma~\ref{lem:sigma-inducing_path_into}, these $\sigma$-inducing paths can be chosen to be into $i$ if the edge is $i \hus j$, into $j$ if the edge is $i \suh j$, and both into $i$ and $j$ if the edge is $i \huh j$. 
  A collider on $\mu$ then stays a collider on the walk between $u$ and $v$ in $G$ obtained by concatenating all these $\sigma$-inducing paths.
  Since such colliders are in $\Anc_H(\{u,v\})$ by assumption, they are also in $\Anc_G(\{u,v\})$ by Lemma~\ref{lem:dpag_ancestral_relations}.
  Each non-endpoint non-collider on $\mu$ must be unblockable.
  But $H$ has no non-trivial strongly connected components, since we already showed that it must be acyclic.
  Therefore, it cannot contain any non-endpoint non-colliders.
  All intermediate colliders on some $\sigma$-inducing path in $G$ between $i$ and $j$ are ancestor in $G$ of $i$ or $j$, and hence, of $u$ or $v$.
  The walk in $G$ obtained by concatenating all $\sigma$-inducing paths along $\mu$ is therefore actually a $\sigma$-inducing walk in $G$.
  So there must also be a $\sigma$-inducing path in $G$ between $u$ and $v$.
  Hence, $u,v$ must be adjacent in $H$, since $H$ represents $G$.

  No input node in $H$ is adjacent to another input node in $H$, because that would mean that there exists a $\sigma$-inducing path in $G$ between the two input nodes, which is impossible.
\end{proof}

The following result shows that every CDMG can be represented by a CDPAG.
\begin{Prp}\label{prp:from_dmg_to_dpag}
  Let $G$ be a CDMG. There exists a CDPAG, $\CDPAG(G)$, that represents $G$.
\end{Prp}
\begin{proof}
  Denote the input nodes of $G$ as $J$, and the output nodes of $G$ as $V$.
  We will construct a mixed graph $H$ with input nodes $J$ and output nodes $V$ as follows.
  Let two vertices $i,j \in J\cup V$ be adjacent in $H$ if and only if there is a $\sigma$-inducing path between $i,j$ in $G$.
  In that case, orient the edge between $i$ and $j$ in $H$ as follows:
    $$\begin{cases}
      \text{$i \ouo j$} & \text{if $i \in \Anc_G(j)$ and $j \in \Anc_G(i)$}, \\
      \text{$i \tuh  j$} & \text{if $i \in \Anc_G(j)$ and $j \notin \Anc_G(i)$}, \\
      \text{$i \hut  j$} & \text{if $i \notin \Anc_G(j)$ and $j \in \Anc_G(i)$}, \\
      \text{$i \huh j$} & \text{if $i \not\in \Anc_G(j)$ and $j \not\in \Anc_G(i)$.}
    \end{cases}$$
%    \begin{compactitem}
%      \item $u \to v \in \tilde{\C{E}}$ if and only if there is a $\sigma$-inducing path between $u$ and $v$ in $\C{G}$ and $u \in \ansub{\C{G}}{v}$,
%      \item $u \ot v \in \tilde{\C{E}}$ if and only if there is a $\sigma$-inducing path between $u$ and $v$ in $\C{G}$ and $v \in \ansub{\C{G}}{u}$,
%      \item $u \oto v \in \tilde{\C{F}}$ if and only if there is a $\sigma$-inducing path between $u$ and $v$ in $\C{G}$ and $u \not\in \ansub{\C{G}}{v}$ and $v \not\in \ansub{\C{G}}{u}$.
%    \end{compactitem}
%This construction preserves the (non-)ancestral relations as well as the $id$-separations/connections.
  It is obvious by construction that $H$ represents $G$.
  It is a valid CDPAG by Proposition~\ref{prp:dpag_represents_dmg_is_valid}.
\end{proof}
If $G$ is acyclic and contains no input nodes, the CDPAG constructed in this way contains no circle edge marks (and is called the directed maximal ancestral graph (DMAG) induced by $G$ in the literature).

Lemma~\ref{lem:ancestral_relations_open_path} can be generalized to CDPAGs:
\begin{Lem}\label{lem:dpag_ancestral_relations_definitely_open_path}
  Let $H$ be a CDPAG that represents some CDMG, with input nodes $J$ and output nodes $V$.
  Let $i, j \in J \cup V$ and $Z \subseteq J \cup V$.
  If $\pi$ is a definitely $Z$-open walk between $i$ and $j$ in $H$, then every node on $\pi$ is in $\Anc_H(\{i,j\} \cup Z)$.
\end{Lem}
\begin{proof}
  Suppose $k$ is a node on $\pi$. 
  Then either $k$ is a definite collider, in which case $k \in \Anc_H(Z)$,
  or there is a directed subwalk from $k$ to a collider on $\pi$, in which case $k \in \Anc_H(Z)$ as well,
  or there is a directed subwalk from $k$ to an endnode of $\pi$, in which case $k \in \Anc_H(\{i,j\})$.
  In all cases, $k \in \Anc_H(\{i,j\}\cup Z)$.
\end{proof}

The following important result states that the existence of a definitely open path in a CDPAG representing a CDMG implies the existence of a corresponding open path in the CDMG.
\begin{Prp}\label{prp:dpag_open_walk_implies_dmg_open_walk}
  Let $H$ be a CDPAG that represents CDMG $G$, both with input nodes $J$ and output nodes $V$.
  Let
    \[\pi = (q_0, a_1, q_1, \dots q_{n-1}, a_n, q_n)\] 
  be a definitely $Z$-open path in $H$, for $Z \subseteq (J \cup V) \sm \{q_0,q_n\}$.
  Then there exists a $Z$-$\sigma$-open path in $G$ between $q_0$ and $q_n$.
\end{Prp}
\begin{proof}
  The vertices $q_0, q_1, \dots, q_n$ in $V \cup J$ must all be distinct.
  For $i=1,\dots,n$, let $\mu_i$ be a $\sigma$-inducing path in $G$ between $q_{i-1}$ and $q_i$ that is into $q_{i-1}$ if the edge $a_{i}$ is into $q_{i-1}$, and into $q_{i}$ if the edge $a_i$ is into $q_i$ (see Lemma~\ref{lem:sigma-inducing_path_into}).
  These can be concatenated into a walk $\mu = (\mu_1,\dots,\mu_n)$ between $q_0$ and $q_n$:
  \[\overunderbraces{&\br{2}{\mu_1}&&&&\br{2}{\mu_{n}}}{&q_0 \cdots \sus & q_1 & \sus \dots \sus q_2 & \sus \cdots \cdots \sus & q_{n-2} \sus \dots \sus & q_{n-1} & \sus \cdots q_n}{&&\br{2}{\mu_2}&&\br{2}{\mu_{n-1}}}\]
  Let $Z \subseteq (J \cup V) \sm \{q_0,q_n\}$.
  
  Consider all walks in $G$ between $q_0$ and $q_n$ with the property that 
  \begin{enumerate}
    \item all colliders on it are in $\Anc_{G}(\{q_0,q_n\} \cup Z)$, and 
    \item each non-endpoint non-collider on it is not in $\{q_0,q_n\} \cup Z$ or is unblockable.
  \end{enumerate}
  Such walks exist, since the concatenation $\mu = (\mu_1, \dots, \mu_n)$ is one, as we will now show. 

  To show the first property, note that by Lemma~\ref{lem:dpag_ancestral_relations_definitely_open_path}, all $q_i \in \Anc_H(\{q_0,q_n\} \cup Z)$, and therefore, by Lemma~\ref{lem:dpag_ancestral_relations}, $q_i \in \Anc_G(\{q_0,q_n\} \cup Z)$.
  This holds in particular for all $q_i$ on $\pi$ that are a collider on $\mu$.
  Every internal collider on some $\mu_i$ is in $\Anc_G(\{q_{i-1},q_i\})$, and hence also these internal colliders are in $\Anc_G(\{q_0,q_n\} \cup Z)$.

  For the second property, note that all internal non-colliders on some $\mu_i$ are unblockable by assumption.
  Suppose that $q_i$ ($1 \le i < n$) is a non-collider on $\mu$. 
  Then $q_i \notin \{q_0,q_n\}$ because all nodes $q_0,\dots,q_n$ are distinct.
  If $q_i$ were in $Z$, then $q_i$ had to be a definite collider on $\pi$, and hence a collider on $\mu$ by construction, contradicting the assumption.
  Hence it is not in $Z$. 
  Therefore, $q_i \notin \{q_0,q_n\} \cup Z$.

  Let $\nu$ be such a walk with a minimal number of colliders.
  We show that all colliders on $\nu$ must be in $\Anc_{G}(Z)$. 

  Suppose on the contrary the existence of a collider $k$ on $\nu$ that is not ancestor of $Z$. 
  It is ancestor of $q_0$ or of $q_n$, by assumption.
  Then there exists a directed path from $k$ to $q_0$ that does not pass through $q_n$, or there exists a directed path from $k$ to $q_n$ that does not pass through $q_0$.
  Without loss of generality, assume the former:
  there is a directed path from $k$ to $q_0$ in $G$ that does not pass through $Z \cup \{q_n\}$.
  Let $\pi$ be a shortest path (possibly trivial) of that type.
  The directed path $\pi$ from $k$ to $q_0$ can be concatenated with the subwalk of $\nu$ between $k$ and $q_n$ (which is into $k$) into a walk between $q_0$ and $q_n$.
  This walk has the property, but has fewer colliders than $\nu$: a contradiction.

%  Therefore, $\nu$ is $\sigma$-open given $Z$ and all occurrences of $q_0$ and $q_n$ as non-endpoint non-collider are unblockable.
  Therefore, $\nu$ is a $Z$-$\sigma$-open walk in $G$ between $q_0$ and $q_n$.
  This means that there must also exist a $Z$-$\sigma$-open path in $G$ between $q_0$ and $q_n$.
\end{proof}

\subsection{Unshielded triples}

One of the key steps in the CFCI algorithm is the orientation of ``unshielded triples''.
The following proposition will later be used to ``orient'' the edges in unshielded triples in a CDPAG representing a CDMG.
\begin{Prp}\label{prp:orienting_unshielded_triples}
  Let $H$ be a CDPAG that represents a CDMG $G$ with input nodes $J$ and output nodes $V$.
  If $\lt i, j, k \rt$ form an unshielded triple in $H$ with $i \in V$ (and $j,k \in V \cup J$), then either 
  \begin{enumerate}[label=(\roman*)]
    \item $j \notin Z$ for each $Z \subseteq (V \cup J) \sm \{i,k\}$ that $i\sigma$-separates $i$ from $k$ in $G$, and $j \notin \Anc_G(\{i,k\})$, or
    \item $j \in Z$ for each $Z \subseteq (V \cup J) \sm \{i,k\}$ that $i\sigma$-separates $i$ from $k$ in $G$, and $j \in \Anc_G(\{i,k\})$.
  \end{enumerate}
\end{Prp}
\begin{proof}
%  Case (i): $j$ is a collider in $H$ on the path $i \sus j \sus k$.
%  Let $Z \subseteq (V \cup J) \sm \{i,k\}$ such that $i \isPerp_G k \given Z$. 
%  If $j \in Z$, the path $i \sus j \sus k$ is $Z$-open in $H$.
%  By Proposition~\ref{prp:dpag_open_walk_implies_dmg_open_walk}, this implies the existence of a $Z$-$\sigma$-open walk in $G$ between $i$ and $k$.
%  Contradiction. 
%  Hence $j \notin Z$.
%  Since $H$ represents $G$, the arrowheads next to $j$ on the path $i \suh j \hus k$ in $H$ imply that $j \notin \Anc_G(\{i,k\})$.
%
%  Case (ii): $j$ is a non-collider in $H$ on the path $i \sus j \sus k$.
%  Let $Z \subseteq (V \cup J) \sm \{i,k\}$ such that $i \isPerp_G k \given Z$. 
%  If $j \notin Z$, the path $i \sus j \sus k$ is $Z$-open in $H$.
%  By Proposition~\ref{prp:dpag_open_walk_implies_dmg_open_walk}, this implies the existence of a $Z$-$\sigma$-open walk in $G$ between $i$ and $k$.
%  Contradiction. 
%  Hence $j \in Z$.
%  Since $H$ represents $G$, the tail(s) next to $j$ on the path $i \sus j \sus k$ in $H$ imply that $j \in \Anc_G(i)$ or $j \in \Anc_G(k)$.

  Case (i): $j \notin \Anc_G(\{i,k\})$. 
  By orienting the edge marks at $j$ on the path $i \sus j \sus k$ in $H$ into $i \suh j \hus k$ (if not already oriented that way), we obtain a CDPAG $\tilde{H}$ that represents $G$.
  Let $Z \subseteq (V \cup J) \sm \{i,k\}$ such that $i \isPerp_G k \given Z$. 
  If $j \in Z$, the path $i \suh j \hus k$ is definitely $Z$-open in $\tilde{H}$.
  By Proposition~\ref{prp:dpag_open_walk_implies_dmg_open_walk}, this implies the existence of a $Z$-$\sigma$-open walk in $G$ between $i$ and $k$, a contradiction.
  Hence $j \notin Z$.

  Case (ii): $j \in \Anc_G(\{i,k\})$.
  By orienting the edge marks at $j$ on the path $i \sus j \sus k$ in $H$ into $i \hut j \sus k$ (if $j \in \Anc_G(i)$) or $i \sus j \tuh k$ (if $j \in \Anc_G(k)$) or $i \hut j \tuh k$ (if $j \in \Anc_G(i) \cap \Anc_G(k)$), we obtain a CDPAG $\tilde{H}$ that represents $G$.
  Let $Z \subseteq (V \cup J) \sm \{i,k\}$ such that $i \isPerp_G k \given Z$. 
  If $j \notin Z$, the path $i \sus j \sus k$ is definitely $Z$-open in $\tilde{H}$.
  By Proposition~\ref{prp:dpag_open_walk_implies_dmg_open_walk}, this implies the existence of a $Z$-$\sigma$-open walk in $G$ between $i$ and $k$, a contradiction.
  Hence $j \in Z$.
\end{proof}
Note that in the first case, we can orient the edges as $i \suh j \hus k$ (if they were not already oriented in this way) to obtain a CDPAG $\tilde{H}$ that also represents $G$. 
In the second case, we cannot orient the edges, since we don't know whether $j \in \Anc_G(i)$ or $j \in \Anc_G(k)$ or both.

\subsection{Discriminating paths}

Another step in FCI is related to the notion of ``discriminating paths''.
This can be considered as an extension of the notion of unshielded triple.
\begin{Def}\label{def:discriminating_path}
  A path $\pi = \lt i, q_1, \dots, q_n, j, k\rt$ (with $n \ge 1$) in a mixed graph $H$ is a \emph{discriminating path for $j$} if:
  \begin{enumerate}[label=(\roman*)]
    \item $i$ is not adjacent to $k$ in $H$, and
    \item every node $q_i$ ($1 \le i \le n$) is a collider on $\pi$ and a parent of $k$ in $H$.
  \end{enumerate}
\end{Def}
Figure~\ref{fig:discriminating_path} illustrates this notion.

\begin{figure}\centering
  \begin{tikzpicture}
    \begin{scope}
      \node[ndout] (X) at (0,0) {$i$};
      \node[ndout] (Q1) at (1.5,0) {$q_1$};
      \node[ndout] (Q2) at (3,0) {$q_2$};
      \node (dots) at (4.5,0) {$\cdots$};
      \node[ndout] (Qn) at (6,0) {$q_n$};
%      \node[ndout] (B) at (7.5,0.0) {$b$};
%      \node[ndout] (Y) at (7.5,2) {$y$};
      \node[ndout] (B) at (7.5,0) {$j$};
      \node[ndout] (Y) at (9,0) {$k$};
      \draw[suh] (X) -- (Q1);
      \draw[arlat] (Q1) -- (Q2);
      \draw[arlat] (Q2) -- (dots);
      \draw[arlat] (dots) -- (Qn);
      \draw[arout,bend right,in=240] (Q1) edge (Y);
      \draw[arout,bend right,in=225] (Q2) edge (Y);
      \draw[arout,bend right,in=210] (Qn) edge (Y);
      \draw[suh]  (B) -- (Qn);
      \draw[sus] (B) -- (Y);
    \end{scope}
  \end{tikzpicture}
  \caption{Discriminating path $\lt i,q_1,q_2,\dots,q_n,j,k \rt$ for $j$ between $i$ and $k$.
  Only $i$ and $j$ could be an input node, all other nodes must be output nodes.\label{fig:discriminating_path}}
\end{figure}

\begin{Rem}\label{rem:discriminating_path}
  It is instructive to think about a discriminating path rather as a certain collection of paths:
  \begin{align*}
    &i \suh q_1 \tuh k \\
    &i \suh q_1 \huh q_2 \tuh k \\
    &\vdots \\
    &i \suh q_1 \huh q_2 \huh \dots \huh q_n \tuh k \\
    &i \suh q_1 \huh q_2 \huh \dots \huh q_n \hus j \sus k
  \end{align*}
  with the additional requirement that $i$ and $k$ are not adjacent.
\end{Rem}

The following quintessential property of discriminating paths is analogous to that of unshielded triples. 
\begin{Prp}\label{prp:orienting_discriminating_paths}
  Let $H$ be a CDPAG that represents a CDMG $G$ with input nodes $J$ and output nodes $V$.
  If $\lt i, q_1, \dots, q_n, j, k \rt$ is a discriminating path in $H$ for $j$ between $i$ and $k$, then
  either 
  \begin{enumerate}[label=(\roman*)]
    \item $j \notin Z$ for each $Z \subseteq (V \cup J) \sm \{i,k\}$ that $i\sigma$-separates $i$ from $k$ in $G$, and $j \notin \Anc_G(\{q_n,k\})$, or
    \item $j \in Z$ for each $Z \subseteq (V \cup J) \sm \{i,k\}$ that $i\sigma$-separates $i$ from $k$ in $G$, and $j \in \Anc_G(k)$.
  \end{enumerate}
  In both cases, $k \notin \Anc_G(j)$.
\end{Prp}
\begin{proof}
  Since $q_n$ is a parent of $k$ in $H$, $q_n \in \Anc_G(k) \sm \{k\}$, and thus $k \in V$.
  We first show that if $Z \subseteq (V \cup J) \sm \{i,k\}$ such that $k \isPerp_G i \given Z$, then $\{q_1, \dots, q_n\} \in Z$.
  This can be seen from the various subpaths in Remark~\ref{rem:discriminating_path}.
  From the first path: if $q_1 \notin Z$ then this path would be definitely open in $H$, and hence (by Proposition~\ref{prp:dpag_open_walk_implies_dmg_open_walk}) there must be a $Z$-$\sigma$-open walk in $G$ between $i$ and $k$, which would contradict $k \isPerp_G i \given Z$.
  Once we have shown that $\{q_1, \dots, q_{i-1}\} \in Z$ for $i \le n$, we see from the $i$'th path that $q_i \in Z$ to avoid definitely opening up this path in $H$, and thereby avoiding the existence of a $Z$-$\sigma$-open path in $G$ between $i$ and $k$.

  Case (i): $j \notin \Anc_G(\{q_n,k\})$.
  By orienting the edge marks at $j$ on the path $q_n \sus j \sus k$ in $H$ into $q_n \suh j \hus k$ (if not already oriented that way), we obtain a CDPAG $\tilde{H}$ that represents $G$.
%  Case (i): $j$ is a collider on the discriminating path.
  Let $Z \subseteq (V \cup J) \sm \{i,k\}$ such that $k \isPerp_G i \given Z$. 
  If $j \in Z$, the discriminating path in $\tilde{H}$ is definitely $Z$-open.
  By Proposition~\ref{prp:dpag_open_walk_implies_dmg_open_walk}, this implies the existence of a $Z$-$\sigma$-open walk in $G$ between $i$ and $k$.
  Contradiction. 
  Hence $j \notin Z$.
  Since $H$ represents $G$, the arrowheads next to $j$ on the subpath $q_n \suh j \hus k$ in $H$ imply that $j \notin \Anc_G(\{q_n,k\})$.

  Case (ii): $j \in \Anc_G(\{q_n,k\})$. 
  Since $q_n \in \Anc_G(k)$, this implies $j \in \Anc_G(k)$.
  By orienting the edge mark at $j$ on the edge $j \sus k$ in $H$ into $j \tuh k$, we obtain a CDPAG $\tilde{H}$ that represents $G$.
%  Case (ii): $j$ is a non-collider in $H$ on the path $i \sus j \sus k$.
  Let $Z \subseteq (V \cup J) \sm \{i,k\}$ such that $k \isPerp_G i \given Z$. 
  If $j \notin Z$, the discriminating path in $\tilde{H}$ is definitely $Z$-open.
  By Proposition~\ref{prp:dpag_open_walk_implies_dmg_open_walk}, this implies the existence of a $Z$-$\sigma$-open walk in $G$ between $i$ and $k$.
  Contradiction. 
  Hence $j \in Z$.

  For the last statement: if $k$ were in $\Anc_G(j)$, then $q_n \in \Anc_G(j)$ because $q_n \in \Anc_G(k)$, contradicting the arrowhead at $q_n$ in $q_n \hus j$.
\end{proof}
Note that in the first case, we can orient $q_n \huh j \huh k$ (as far as the edge marks were not already oriented in this way) to obtain a CDPAG $\tilde{H}$ that also represents $G$. 
In the second case, we can orient $j \tuh k$ (as far as the edge marks were not already oriented in this way) to obtain a CDPAG $\tilde{H}$ that also represents $G$.

\subsection{Conditional independence models and Markov equivalence}

\begin{Def}\label{def:independence_model}
  For a CDMG $G$ with input nodes $J$ and output nodes $V$, define its \emph{conditional $d$-independence model} to be
  $$\IM_d(G) := \{ \lt A, B, C \rt : A, B, C \subseteq J \cup V, A \cap J = \emptyset, J \subseteq B \cup C : A \idPerp_G B \mid C \},$$
  i.e., the set of all $id$-separations of restricted form entailed by the graph. Define its
  \emph{conditional $\sigma$-independence model} to be
  $$\IM_\sigma(G) := \{ \lt A, B, C \rt : A, B, C \subseteq J \cup V, A \cap J = \emptyset, J \subseteq B \cup C : A \isPerp_G B \mid C \},$$
  i.e., the set of all $i\sigma$-separations of restricted form entailed by the graph. 
\end{Def}
In general, given two disjoint index sets $J,V$, we call a subset of 
$\{ \lt A, B, C \rt : A, B, C \subseteq J \cup V, A \cap J = \emptyset, J \subseteq B \cup C  \}$
a \emph{conditional independence model over $V \given J$}.
If $\lt A, B, C \rt$ is an element of a conditional independence model, we also say that $C$ separates $A$ from $B$.
Both $\IM_d(G)$ and $\IM_\sigma(G)$ are conditional independence models over $V \given J$ (with
$J$ the input nodes of $G$ and $V$ the output nodes of $G$).
For CADMGs, $i\sigma$-separation is equivalent to $id$-separation, and hence, if $G$ is acyclic, then
$\IM_d(G) = \IM_\sigma(G)$.

To connect with the literature, we also provide the following definition in case no exogenous input nodes are present:
\begin{Def}\label{def:independence_model}
  For a DMG $G$ with vertex set $V$, define its \emph{$d$-independence model} to be
  $$\IM_d(G) := \{ \lt A, B, C \rt : A, B, C \subseteq V, A \dPerp_G B \mid C \},$$
  i.e., the set of all $id$-separations entailed by the graph. Define its
  \emph{$\sigma$-independence model} to be
  $$\IM_\sigma(G) := \{ \lt A, B, C \rt : A, B, C \subseteq V, A \sPerp_G B \mid C \},$$
  i.e., the set of all $i\sigma$-separations entailed by the graph. 
\end{Def}
In general, given an index set $V$, we call a subset of 
$\{ \lt A, B, C \rt : A, B, C \subseteq V \}$
an \emph{independence model over $V$}.
If $\lt A, B, C \rt$ in an independence model, we also say that $C$ separates $A$ from $B$.
Both $\IM_d(G)$ and $\IM_\sigma(G)$ are independence models over $V$, the vertices of $G$.
For ADMGs, $i\sigma$-separation is equivalent to $id$-separation, and hence, if $G$ is acyclic, then
$\IM_d(G) = \IM_\sigma(G)$.

The input to the original FCI algorithm consists of an independence model over $V$,
and its output will consist of a mixed graph with vertices $V$.
If the input of FCI is the independence model of a DMG, then the output of FCI
will be a DPAG that represents the ``Markov equivalence class'' of the DMG.
\begin{Def}\label{def:Markov_equivalent}
  We call two DMGs $G_1$ and $G_2$ \emph{$\sigma$-Markov equivalent} if $\IM_\sigma(G_1) = \IM_\sigma(G_2)$,
  and \emph{$d$-Markov equivalent} if $\IM_d(G_1) = \IM_d(G_2)$.
\end{Def}
The input to the CFCI algorithm will consist of a conditional independence model over $V \given J$,
and its output will consist of a conditional mixed graph with input nodes $J$ and output nodes $V$.
If the input of CFCI is the conditional independence model of a CDMG, then the output of FCI
will be a CDPAG that represents the ``Markov equivalence class'' of the CDMG, as
we will see later.
\begin{Def}\label{def:Markov_equivalent}
  We call two CDMGs $G_1$ and $G_2$ \emph{$\sigma$-Markov equivalent} if $\IM_\sigma(G_1) = \IM_\sigma(G_2)$,
  and \emph{$d$-Markov equivalent} if $\IM_d(G_1) = \IM_d(G_2)$.
\end{Def}

\subsection{Skeleton search}

The CFCI algorithm consists of two phases, the skeleton search phase, which is followed by the orientation phase.
In this subsection, we describe the skeleton search phase.

\begin{Def}\label{def:skeleton}
Given a CDPAG $H = \lt J, V, E \rt$, its \emph{skeleton} is the mixed graph 
$\skel(H) = \lt J, V, F \rt$ with the same nodes, and with a bicircle edge $i \ouo j$ in $F$
if and only if $i \sus j$ in $E$ (i.e., if $i$ and $j$ are adjacent in $H$).
\end{Def}
Hence, the only edge type occurring in the skeleton is the bicircle edge.
It has no edge between any pair of nodes in $J$.
We will later frequently refer to the unordered pairs of nodes that may be adjacent:
\begin{Def}
  For input nodes $V$ and output nodes $J$, define $\separable(V,J) := \{ \{i,j\} : i \in V, j \in J \cup V, i \ne j \}$.
\end{Def}

The aim of the skeleton search phase of the CFCI algorithm is to construct the skeleton of the CDPAG that represents a CDMG from the conditional $\sigma$-independence model of the CDMG.
It does this by testing for each edge that might be absent in the skeleton whether it can find any subset that separates the two nodes.
If it finds a separating set between an (unordered) pair of distinct nodes $\{i, j\}$, the set is stored in $\sepset(\{i,j\})$.
The orientation phase later makes use of these separating sets found in the skeleton phase.

A brute-force search over all possible subsets of $(J \cup V)\sm \{i,j\}$, as in Algorithm~\ref{alg:bf_skeleton}, would be a straightforward solution.
However, it is also computationally extremely expensive for all but the smallest cardinalities of $V$ and $J$, and statistically not very reliable in case the separations have to be tested with conditional independence tests on finite data.
\begin{algorithm}
  \begin{algorithmic}[1]
    \Input Input node index set $J$; output node index set $V$; conditional independence model $I$ over $V \given J$
    \Output mixed graph $H$ with input nodes $J$ and output nodes $V$; separating sets $\sepset$
    \State $H \leftarrow \lt J, V, \emptyset \rt$
    \ForEach{$\{i,j\} \in \separable(V,J)$}
      \State add edge $i \ouo j$ to $H$
      \ForEach{$Z \subseteq (V \cup J) \sm \{i,j\}$}
      \If{$\lt i, j, Z \rt \in I$}\Comment{found a separating set}
        \State delete edge $i \ouo j$ from $H$
        \State $\sepset(\{i,j\}) \leftarrow Z$
        \Break
      \EndIf
      \EndFor
    \EndFor
  \end{algorithmic}
  \caption{Brute-force skeleton algorithm $\mathtt{BFskeleton}(J,V,I)$.\label{alg:bf_skeleton}}
\end{algorithm}

To get some inspiration on how to address this, we will first describe the skeleton search phase of the PC algorithm, an ancestor of the FCI algorithm designed for DAGs \cite{SGS2000}.
The PC skeleton phase (Algorithm~\ref{alg:pc_skeleton}) searches for separating sets of increasing cardinality. 
As candidates for a separating set between nodes $i,j \in H$, it considers all subsets of $\Adj_H(i)$ and all subsets of $\Adj_H(j)$.
The rationale is that if $G$ is a DAG, then either $\Pa_G(i) \subseteq \Adj_H(i)$ or $\Pa_G(j) \subseteq \Adj_H(j)$ $id$-separates $i$ from $j$.

\begin{algorithm}
  \begin{algorithmic}[1]
  \Input independence model $I$ over vertex set $V$
  \Output mixed graph $H$ with vertex set $V$; separating sets $\sepset$
  \State initialize $H$ as a complete graph with only bicircle ($\ouo$) edges
  \State $n \leftarrow 0$
    \Repeat
    \Repeat
    \State select ordered pair $i \ouo j$ in $H$ such that $\#(\Adj_H(i)\sm\{j\}) \ge n$
    \State select a subset $Z \subseteq \Adj_H(i)\sm\{j\}$ of cardinality $n$
    \If{$(i,j,Z) \in I$} 
      \State delete edge $i \ouo j$ from $H$
      \State $\sepset(\{i,j\}) \leftarrow Z$
    \EndIf
    \Until{no more such tuples $(i,j,Z)$ can be selected}
    \State $n \leftarrow n + 1$
    \Until{for each ordered pair $i \sus j$ in $H$, $\#(\Adj_H(i)\sm\{j\}) < n$}
  \end{algorithmic}
  \caption{PC skeleton algorithm $\mathtt{PCskeleton}(V,I)$. \Joris{remove?}\label{alg:pc_skeleton}}
\end{algorithm}

We can easily extend this to deal with input nodes as well, see Algorithm~\ref{alg:pc_skeleton}.
Since the input is a \emph{conditional} independence model, all separating sets need to contain all nodes in $J$ (except $j$ itself, if $j \in J$).
We therefore only need to consider subsets of neighbours of $i$ that are not in $J$, in increasing cardinality.

\begin{algorithm}
  \begin{algorithmic}[1]
  \Input Input node index set $J$; output node index set $V$; conditional independence model $I$ over $V \given J$
  \Output mixed graph $H$ with input nodes $J$ and output nodes $V$; separating sets $\sepset$
  \State $H \leftarrow \lt J, V, \emptyset \rt$
  \ForEach{$\{i,j\} \in \separable(V,J)$}
    \State add edge $i \ouo j$ to $H$
  \EndFor
  \State $n \leftarrow 0$
    \Repeat
    \Repeat
    \State select $i \ouo j$ in $H$ with $\#(\Adj_H(i)\sm(J\cup\{j\})) \ge n$
    \State select a subset $Z \subseteq \Adj_H(i)\sm(J\cup\{j\})$ of cardinality $n$
    \If{$(i,j,Z \cup (J\sm\{j\})) \in I$} 
      \State delete edge $i \ouo j$ from $H$
      \State $\sepset(\{i,j\}) \leftarrow Z \cup (J \sm \{j\})$
    \EndIf
    \Until{no more such tuples $(i,j,Z)$ can be selected}
    \State $n \leftarrow n + 1$
    \Until{for each $i \ouo j$ in $H$, $\#(\Adj_H(i)\sm(J\cup\{j\})) < n$}
  \end{algorithmic}
  \caption{Conditional PC skeleton algorithm $\mathtt{CPCskeleton}(J,V,I)$.\label{alg:pc_skeleton}}
\end{algorithm}

The underlying idea may no longer hold if $G$ is an ADMG or a DMG, or in case of selection bias.
The original proposal (which motivated the somewhat optimistic adjective ``Fast'' in the name of the FCI algorithm \cite{SMR1995}) replaces the subsets of adjacent nodes by so-called ``Possible-D-Sep'' sets. 
Before we can define these, we need some definitions and theory.

\begin{Def}\label{def:SEP}
  Let $G$ be a CDMG with input nodes $J$ and output nodes $V$.
  For $i,j \in J \cup V$ distinct, define $\SEP_G(i,j)$ to be the set of nodes $k \in V$ such that
  $k \ne i$ and there is a walk $\pi$ in $G$ between $i$ and $k$ such that every node
  on $\pi$ is in $\Anc_G(\{i,j\})$, and every non-endpoint non-collider on $\pi$ is unblockable.
\end{Def}
The name $\SEP_G(i,j)$ is motivated by the following property.
\begin{Prp}\label{prp:SEP}
  Let $G$ be a CDMG with input nodes $J$ and output nodes $V$.
  Let $i \in V$ and $j \in J \cup V$ be distinct.
  If there exists no $\sigma$-inducing walk between $i,j$ in $G$, then $\SEP_G(i,j) \cap \{i,j\} = \emptyset$ and $i \isPerp_G j \given \SEP_G(i,j) \cup (J \sm \{j\})$.
\end{Prp}
\begin{proof}
  By definition, $\SEP_G(i,j) \subseteq \Anc_G(\{i,j\})$ and $i \notin \SEP_G(i,j)$.
  If $j \in \SEP_G(i,j)$, the walk $\pi$ in Definition~\ref{def:SEP} between $i$ and $j$ must be $\sigma$-inducing.
  Since there exists no $\sigma$-inducing walk between $i,j$ in $G$ by assumption, $j \notin \SEP_G(i,j)$.

  We prove the $i\sigma$-separation by contradiction.
  Write $Z := \SEP_G(i,j) \cup (J \sm \{j\})$.
  Suppose there exists a path in $G$ between $i$ and $\{j\} \cup J$ that is $\sigma$-open given $Z$.
  It cannot end in a node in $J \sm \{j\}$, and hence it must end in $j$.
  Let $\pi = v_0 \sus \dots \sus v_n$ with $v_0 = i$ and $v_n \in j$ be such a path consisting of a minimal number of nodes.
  Every node on $\pi$ must be in $\Anc_G(\{i,j\}) \cup J$, since by Lemma~\ref{lem:ancestral_relations_open_path}, each node on $\pi$ is in $\Anc_G(\{i,j\} \cup Z)$, and $Z \subseteq \Anc_G(\{i,j\}) \cup J$.
  $\pi$ can only contain nodes in $J$ as endnodes (otherwise we could shorten the path).
  In other words, the only node on $\pi$ that could be in $J$ is $j$.

  Denote the subwalk of $\pi$ from $v_a$ to (and including) $v_b$ by $\pi(v_a,v_b)$, for $0 \le a \le b \le n$.
  We will show that for all $k = 1,\dots,n$, $\pi(v_0,v_k)$ has the property that every non-endpoint non-collider on it is unblockable.
  The property trivially holds for $k=1$. 
  Suppose it holds for $k < n$.
  Since $v_0,\dots,v_k$ are all in $\Anc_G(\{i,j\})$, and all non-endpoint non-colliders on $\pi(v_0,v_k)$ are unblockable, we conclude that $v_k \in \SEP_G(i,j)$.
  If $v_{k}$ is a non-collider on $\pi(v_0,v_{k+1})$, it must be unblockable, because $\pi(v_0,v_{k+1})$ is $Z$-$\sigma$-open and $v_k \in \SEP_G(i,j) \subseteq Z$.
  So the property also holds for $k+1$.

  In particular, we can conclude that $j = v_n \in \SEP_G(i,j)$, a contradiction.
\end{proof}
In practice, one does not know the set $\SEP_G(i,j)$ if $G$ is unknown.
One can, however, easily obtain a `bound' on this set by identifying a superset.
\begin{Def}\label{def:possep}
  Let $H_0 = \lt J,V,E\rt$ be a mixed graph with input nodes $J$, output nodes $V$ and edges $E$ of the types $\{\tuh,\hut,\huo,\huh,\ouo,\ouh\}$, with at most one edge connecting any pair of distinct nodes.
  For $i \in V, j \in J \cup V$ distinct, we define $\PosSEP_{H_0}(i,j) \subseteq V$ to consist of those nodes
  $k \in V$ such that $k \not\in \{i,j\}$ and there is a path between $i$ and $k$ in $H_0$ such that for
  every subsequent triple $a \sus b \sus c$ on the path, either the triple is a definite collider in $H_0$,
  or a triangle in $H_0$.
\end{Def}
We can now show that:
\begin{Lem}\label{lem:possep}
  Let $G$ be a CDMG with input nodes $J$ and output nodes $V$.
  Let $H_0$ be the mixed graph constructed by lines 
  \ref{alg:fci_skeleton2:cpc_skeleton}--\ref{alg:fci_skeleton2:oriented_unshielded_colliders}
  of Algorithm~\ref{alg:fci_skeleton} when run on the conditional independence model $\IM_\sigma(G)$.
  Then, for $i \in V, j \in J \cup V$ distinct, if there is no $\sigma$-inducing path in $G$ between $i$ and $j$, then
  $\SEP_G(i,j) \subseteq \PosSEP_{H_0}(i,j)$.
\end{Lem}
\begin{proof}
%  That is, first one runs the Conditional PC skeleton algorithm (Algorithm~\ref{alg:pc_skeleton}),
%  then one orients all unshielded colliders as definite colliders if the middle node is in the separating set of the two outer mnodes, and finally one orients the edges between input and output nodes.
  Note that the skeleton of $H_0$ is a supergraph of the skeleton of $\CDPAG(G)$, i.e., every edge in $\CDPAG(G)$ must also be present in $H_0$.
  Let $k \in \SEP_G(i,j)$.
  Since there is no $\sigma$-inducing path in $G$ between $i$ and $j$, $\{i,j\} \cap \SEP_G(i,j) = \emptyset$ by Proposition~\ref{prp:SEP}, and hence $k \ne i$ and $k \ne j$.
  By definition, there is a path $\pi$ in $G$ between $i$ and $k$ with the property that every node on $\pi$ is in $\Anc_G(\{i,j\})$, and every non-endpoint non-collider on $\pi$ is unblockable, and every node on $\pi$ is in $V$.
  There is a corresponding path $\pi'$ in $H_0$ between $i$ and $k$ that consists of the same sequence of nodes (but may have different edges between the nodes).
  Consider a subsequent triple $a \sus b \sus c$ on $\pi$.
  Suppose $b$ is a collider on $\pi$.
  In $H_0$, $a$ must also be adjacent to $b$, and $b$ to $c$.
  If $a$ and $c$ are not adjacent in $H_0$, then a separating set for $a$ and $c$ was found during the construction of $H_0$, or $a$ and $c$ are both in $J$.
  The latter would be a contradiction. 
  Therefore, it will have been oriented as a definite collider in $H_0$ according to Proposition~\ref{prp:orienting_unshielded_triples}.
  Otherwise (if $a$ and $c$ are adjacent in $H_0$), it forms a triangle in $H_0$.
  If $b$ is a non-collider on $\pi$, then it must be unblockable.
  But in that case, $a \sus b \sus c$ is a $\sigma$-inducing walk between $a$ and $c$ in $G$.
  Hence, $(a,b,c)$ will form a triangle in $H_0$.
  Therefore, $k \in \PosSEP_{H_0}(i,j)$.
\end{proof}

\begin{algorithm}
  \begin{algorithmic}[1]
  \Input Input node index set $J$; output node index set $V$; conditional independence model $I$ over $V \given J$
  \Output mixed graph $H$ with input nodes $J$ and output nodes $V$; separating sets $\sepset$
  \State $\lt H_0,\sepset \rt \leftarrow \mathtt{CPCskeleton}(J,V,I)$ \label{alg:fci_skeleton2:cpc_skeleton}
%  \For{$j \ouo v$ in $H$ with $j \in J$, $v \in V$} \Comment{apply background knowledge \Joris{REDUNDANT?}}
%    \State orient the edge as $j \tuh v$
%  \EndFor
%  \State Orient unshielded triples in $H_0$:
%    \begin{description}
%      \item[$\C{R}0$] for each unshielded triple $\lt i,j,k \rt$ in $H_0$ with $j,k \in V$, orient it as $i \suh j \hus k$ if $j \notin \sepset(k,i)$
%    \end{description}\label{alg:fci_skeleton2:oriented_unshielded_colliders}
  \ForEach{unshielded triple $\lt i,j,k \rt$ in $H_0$ with $j,k \in V$} \Comment{orient colliders}
    \If{$j \notin \sepset(\{i,k\})$} % \sepset(k,i) if ordered!
      \State orient it as $i \suh j \hus k$
    \EndIf
  \EndFor\label{alg:fci_skeleton2:oriented_unshielded_colliders}
%  \For{$i \sus j \sus k$ in $H_0$ with $j,k \in V$}  \Comment{orient unshielded colliders}
%    \If{$i \notin \Adj_{H_0}(k)$}
%      \If{$j \notin \sepset(k,i)$}
%        \State orient it as $i \suh j \hus k$
%      \EndIf
%    \EndIf
%  \EndFor \label{alg:fci_skeleton2:oriented_unshielded_colliders}
  %find PossibleDSep sets
  \ForEach{$i \in V, j \in V \cup J$ with $i \ne j$} \Comment{construct PossibleDSep sets}
    \State calculate $\PosSEP_{H_0}(i,j)$
  \EndFor
  %remove orientations
    \State $H \leftarrow (J,V,\{i \ouo j : i \sus j \in H_0 \})$ \Comment{forget orientations} \label{alg:fci_skeleton2:construct_skeleton_start}
%  \For{$(i,j) \in \separable(V,J)$} \Comment{remove orientations}
%    \If{$i \sus j$ in $H_0$}
%      \State replace the edge by $i \ouo j$
%    \EndIf
%  \EndFor
  %find more separating sets
  \State $n \leftarrow 0$
    \Repeat \Comment{search for separating sets}
    \Repeat
    \State select $i \ouo j$ in $H$ with $\#(\PosSEP_{H_0}(i,j)\sm J) \ge n$
    \State select a subset $Z \subseteq \PosSEP_{H_0}(i,j)\sm J$ of cardinality $n$
    \If{$(i,j,Z \cup (J\sm\{j\})) \in I$} 
      \State delete edge $i \ouo j$ from $H$
      \State $\sepset(\{i,j\}) \leftarrow Z \cup (J \sm \{j\})$
    \EndIf
    \Until{no more such tuples $(i,j,Z)$ can be selected}
    \State $n \leftarrow n + 1$
    \Until{for each $i \ouo j$ in $H$, $\#(\PosSEP_{H_0}(i,j)\sm J) < n$}\label{alg:fci_skeleton2:construct_skeleton_end}
  \end{algorithmic}
  \caption{CFCI skeleton algorithm $\mathtt{CFCIskeleton}(J,V,I)$.\label{alg:fci_skeleton}}
\end{algorithm}
%The FCI skeleton search uses the strategy of the CPC skeleton algorithm, but uses $\PosSEP_{H_0}(i,j)$ instead of $\Adj_H(i)\sm\{j\}$ and check for each subset whether it separates $i$ from $j$ (when conditioning also on $J \sm \{j\}$). 
%Lemma~\ref{lem:possep} guarantees that such a separating set exists, if any separating set of $i$ and $j$ exists. 

So we do not need to search over all possible subsets of $(J \cup V) \sm \{i,j\}$ for a separating set between $i,j$, but only the subsets in $\PosSEP_{H_0}(i,j)$.
The skeleton phase of the CFCI algorithm is described in Algorithm~\ref{alg:fci_skeleton}. 
It is an extension of the original FCI skeleton search \cite{SMR1995} to deal with input nodes. 
It first runs the PC skeleton phase and orients triples using $\C{R}0$, obtaining a directed mixed graph $H_0$. 
Then it calculates the sets $\PosSEP_{H_0}(i,j)$ for all distinct pairs $(i,j)$ with $i \in V, j \in J \cup V$.
If $i \isPerp_G j \mid Z$ for some $Z \subseteq (J \cup V) \sm \{i,j\}$, then
$i \isPerp_G j \mid Z^*$ for some $Z^* \subseteq \PosSEP_{H_0}(i,j)$.
Since some of the oriented colliders may be incorrect in $H_0$ (because the PC skeleton phase may have missed some separating sets), it removes all orientations from $H_0$ and then continues with a more extensive search for separating sets, similar
to how the PC skeleton phase is done, but now using $\PosSEP_{H_0}(i,j)$ instead of $\Adj_{H_0}(i)\sm\{j\}$ to find candidate nodes for the separating set.

\begin{Thm}\label{thm:fci_skeleton_sound}
  The CFCIskeleton algorithm (Algorithm~\ref{alg:fci_skeleton}) is \emph{sound}: if its input consists of the conditional $\sigma$-independence model $\IM_\sigma(G)$ of a CDMG $G$, then its output will be $\skel(\CDPAG(G))$.
\end{Thm}
\begin{proof}
  Let $G$ be a CDMG with input nodes $J$ and output nodes $V$ and $I = \IM_\sigma(G)$ its conditional $\sigma$-independence model.
  By Lemma~\ref{lem:possep}, the mixed graph $H_0$ constructed by lines 
  \ref{alg:fci_skeleton2:cpc_skeleton}--\ref{alg:fci_skeleton2:oriented_unshielded_colliders}
  has the property that for distinct $i \in V, j \in J \cup V$, if there is no $\sigma$-inducing path in $G$ between $i$ and $j$, then $\SEP_G(i,j) \subseteq \PosSEP_{H_0}(i,j)$.
  Thus, if $i \in V$ and $j \in J \cup V$ are not adjacent in $\skel(\CDPAG(G))$, then we are guaranteed that for some subset $Z \subseteq \PosSEP_{H_0}(i,j)$, we have that $i \isPerp_G j \given (Z \cup J) \sm \{j\} \cup S$.
%  We have to show that which means that it will compute $H = \skel(\CDPAG(G))$, and $\sepset$ will contain a separating set of $i$ and $j$ for every edge $i \ouo j$ absent in $H$ with $i \in V, j \in J \cup V, i \ne j$.
  The graph $H$ constructed in lines \ref{alg:fci_skeleton2:construct_skeleton_start}--\ref{alg:fci_skeleton2:construct_skeleton_end} will be a mixed graph with input nodes $J$ and output nodes $V$ that only has bicircle edges.
  It has an edge between any pair of distinct nodes $i,j \in J\cup V$ if and only if $\{i,j\} \cap V \ne \emptyset$ and there is no set $Z \subseteq (J \cup V) \sm \{i,j\}$ such that $i \isPerp_G j \given Z \cup S$.
  Furthermore, if there is no edge in $H$ between a pair of distinct nodes $i,j \in J \cup V$, then either $\{i,j\} \subseteq J$, or $i \isPerp_G j \given \sepset(\{i,j\}) \cup S$.
  By Proposition~\ref{prp:sigma-inducing_path_properties}, this implies that two distinct nodes $i,j \in J \cup V$ are adjacent in $H$ at this stage if and only if there is a $\sigma$-inducing walk in $G$ between $i$ and $j$.
  Hence $H$ must be $\skel(\CDPAG(G))$.
\end{proof}

An alternative search strategy was proposed that can be considerably faster in practice, for which it can be shown that the corresponding FCI+ algorithm is of polynomial-time complexity in the number of nodes, as long as the degree of the DPAG is bounded \cite{ClaassenMooijHeskes_UAI_13}.
%Another variant, the ``Really Fast Causal Inference'' algorithm has also been proposed, which adapts both the skeleton phase and the orientation rules, and sacrifices some completeness but remains sound \cite{Colombo++2012}.

\subsection{CFCI Algorithm}

We are now ready to describe a causal inference algorithm, Conditional FCI, that is an extension of the original
Fast Causal Inference (FCI) algorithm of \cite{SMR1995} to deal with input nodes and cycles.
It is presented in Algorithm~\ref{alg:cfci}.
Its input is a conditional independence model over $V \given J$, where $V$ and $J$ are index sets of output and input nodes, respectively. 
Its output is a mixed graph with input nodes $J$ and output nodes $V$.
It starts with a \emph{skeleton phase} (line~\ref{alg:cfci:skeleton}, see Algorithm~\ref{alg:fci_skeleton}) that is aimed at deducing the adjacencies between the nodes, and to find sets that separate two nodes.
Then, it runs various \emph{orientation rules} (lines~\ref{alg:cfci:start_orient}--\ref{alg:cfci:end_orient}) that iteratively orient edge marks into tails and arrowheads.
We have omitted orientation rules $\C{R}5$--$\C{R}7$ of \cite{Zhang2008_AI}, since these are only necessary if selection bias is not ruled out a priori. 
Those rules can yield edges of the form $\tuo, \out$ and $\tut$ that are not valid in CDPAGs, and require the more general
formalism of CPAGs (which is yet to be developed for the general case with possible cycles and input nodes).

Compared to the original FCI algorithm \cite{SMR1995} (which assumes no input nodes), 
we have slightly changed the skeleton search phase (by starting with a mixed graph that
contains no edges between input nodes, limiting the paths in the calculation of the
$\PosSEP_{H_0}(i,j)$ sets to output nodes only, and by including all input nodes, 
except $j$ itself, into the separating set). Furthermore, we added 
step~\ref{alg:cfci:orientJ} to orient the edges between input and output nodes. 
The rest of the algorithm is unchanged.

\begin{algorithm}
\begin{algorithmic}[1]
  \Input Input node index set $J$; output node index set $V$; conditional independence model $I$ over $V \given J$
  \Output mixed graph $H$ with input nodes $J$ and output nodes $V$; separating sets $\sepset$
  \State $\lt H,\sepset \rt \leftarrow \mathtt{CFCIskeleton}(J,V,I)$\label{alg:cfci:skeleton}
  \ForEach{edge $j \ouo v$ in $H$ with $j \in J$, $v \in V$}\label{alg:cfci:orientJ}\label{alg:cfci:start_orient}
    \State orient $j \tuh v$
  \EndFor
%  \For{each edge $j \ouo v$ in $H$ with $j \in J$ and $v \in V$} \label{alg:cfci:orientJ} \Comment{$\C{R}bk$}
%  \State orient it as $j \tuh v$ \Joris{Or $j \ouh v$?} 
%  \EndFor
%  \For{each unshielded triple $\lt i,j,k \rt$ in $H$ with $k \in V$} \Comment{$\C{R}0$}
%    \If{$j \notin \sepset(k,i)$}
%      \State orient it as $i \suh j \hus k$
%    \EndIf
%  \EndFor
  \Repeat \label{alg:fci:R0}
    \begin{description}
      \item[$\C{R}0$] if $i \sus j \sus k$ in $H$, and $i$ and $k$ are not adjacent in $H$, then orient $i \suh j \hus k$ if $j \notin \sepset(\{i,k\})$ % \sepset(k,i) if ordered!
    \end{description}
  \Until{this orientation rule is not applicable}
  \Repeat
    \begin{description}
      \item[$\C{R}1$] if $i \suh j \ous k$ in $H$, and $i$ and $k$ are not adjacent in $H$, then orient $i \suh j \tuh k$
      \item[$\C{R}2$] if $i \tuh j \suh k$ or $i \suh j \tuh k$ in $H$, and $i \suo k$ in $H$, then orient $i \suh k$
      \item[$\C{R}3$] if $i \suh j \hus k$ in $H$, and $i \suo l \ous k$ in $H$, and $l \suo j$ in $H$, and $i$ and $k$ are not adjacent in $H$, then orient $l \suh j$
      \item[$\C{R}4$] if $\lt i, q_1, \dots, q_n, j, k \rt$ is a discriminating path in $H$ for $j$, and if $j \ous k$ in $H$, then orient $j \tuh k$ if $j \in \sepset(\{i,k\})$ and orient $q_n \huh j \huh k$ if $j \notin \sepset(\{i,k\})$ % \sepset(k,i) if ordered!
    \end{description}
  \Until{none of these orientation rules is applicable}
  \Repeat
    \begin{description}
      \item[$\C{R}8$] if $i \tuh j \tuh k$ in $H$, and $i \ouh k$ in $H$, then orient $i \tuh k$
      \item[$\C{R}9$] if $i \ouh k$, and $\pi = \lt i, j, \dots, k \rt$ is an uncovered possibly directed path in $H$ from $i$ to $k$ such that $j$ and $k$ are not adjacent in $H$, then orient $i \tuh k$
      \item[$\C{R}10$] if $i \ouh k$ in $H$, $j \tuh k \hut l$ in $H$, $\pi_1$ is a uncovered possibly directed path in $H$ from $i$ to $j$, and $\pi_2$ is a uncovered possibly directed path in $H$ from $i$ to $l$, then let $u_1$ be the vertex adjacent to $i$ on $\pi_1$ (possibly $u_1 = j$) and $u_2$ the vertex adjacent to $i$ on $\pi_2$ (possible $u_2 = l$); if $u_1 \ne u_2$, and $u_1$ and $u_2$ are not adjacent in $H$, then orient $i \tuh k$
    \end{description}
  \Until{none of these orientation rules is applicable}\label{alg:cfci:end_orient}
  \end{algorithmic}
  \caption{CFCI Algorithm.\label{alg:cfci}}
\end{algorithm}


\begin{Thm}\label{thm:fci_sound}
  The CFCI algorithm (Algorithm~\ref{alg:cfci}) is \emph{sound}: if its input consists of the conditional $\sigma$-independence model $\IM_\sigma(G)$ of a CDMG $G$, then its output will be a valid CDPAG $H$ that represents $G$.
\end{Thm}
\begin{proof}
  Let $G$ be a CDMG with input nodes $J$ and output nodes $V$ and $I = \IM_\sigma(G)$ its conditional $\sigma$-independence model.

  The CFCI skeleton phase in line \ref{alg:cfci:skeleton} is sound (Theorem~\ref{thm:fci_skeleton_sound}),
  that is, it computes $H = \skel(\CDPAG(G))$, and $\sepset$ will contain a separating set of $i$ and $j$ for every edge $i \ouo j$ absent in $H$ with $i \in V, j \in J \cup V, i \ne j$.
  Obviously, the skeleton $\skel(\CDPAG(G))$ is a valid CDPAG that represents $G$.
  The next step, line \ref{alg:cfci:orientJ}, orients the edges between input and output nodes in $H$ as they are in $G$. 
  The result is still a valid CDPAG that represents $G$.
  The rest of the proof proceeds by induction.

  We show for each of the orientation rules, that under the assumption that the current $H$ is a valid CDPAG that represents $G$, applying the rule yields an updated $H$ that is still a valid CDPAG that represents $G$.
  Some of the later rules ($\C{R}1$, $\C{R}3$, $\C{R}9$, $\C{R}10$) assume that rule $\C{R}0$ is no longer applicable, which is the reason that rule $\C{R}0$ is performed before the other orientation rules are performed.

  In the following, we will always assume that the antecedent of the rule holds for a mixed graph $H$ that is a valid CDPAG that represents CDMG $G$.
%  \Joris{Nodes next to a tail cannot be in $J$, and nodes next to a circle cannot be in $J$.}
  \begin{description}
    \item[$\C{R}0$] ``If $i \sus j \sus k$ in $H$, and $i$ and $k$ are not adjacent in $H$, then orient $i \suh j \hus k$ if $j \notin \sepset(\{i,k\})$.''\\ % \sepset(k,i) if ordered!
      It follows from Proposition~\ref{prp:orienting_unshielded_triples} that $H$ still represents $G$ after the orientation of the unshielded collider.
      Note that there is no need to assume $j \in V$, because if $j \in J$, then it will be contained in $\sepset(\{i,k\})$.
    \item[$\C{R}1$] ``If $i \suh j \ous k$ in $H$, and $i$ and $k$ are not adjacent in $H$, orient the triple as $i \suh j \tuh k$.''\\
%      \Joris{Note that $j,k \notin J$, but $i$ could be in $J$. The proof works in both cases and there is no risk of overwriting background knowledge.}
      We have to show that $j \in \Anc_G(k)$.
      Since the triple $\lt i, j, k \rt$ is an unshielded triple in $H$, but has not been oriented as a collider by $\C{R}0$, we conclude that $j \in \sepset(\{i,k\})$. % \sepset(k,i) if ordered!
      By Proposition~\ref{prp:orienting_unshielded_triples}, $j \in \Anc_G(\{i,k\})$.
      Since $H$ represents $G$ and $i \suh j$ in $H$, $j \not\in \Anc_G(i)$. 
      Therefore, $j \in \Anc_G(k)$.
      After orienting $j \tuh k$ in $H$, the resulting mixed graph $H$ still represents $G$.
    \item[$\C{R}2$] ``If $i \tuh j \suh k$ or $i \suh j \tuh k$ in $H$, and $i \suo k$ in $H$, then orient $i \suh k$.''\\
%      \Joris{Note that $j,k \notin J$, but $i$ might be in $J$. There is no risk of overwriting background knowledge.}
      By the antecedent of the rule, and since $H$ represents $G$, we have $k \notin \Anc_G(j)$.
      In case $i \tuh j \suh k$, we have $i \in \Anc_G(j)$, so if $k$ were in $\Anc_G(i)$, it would follow that $k \in \Anc_G(j)$, a contradiction.
      In case $i \suh j \tuh k$, we have on one hand $j \notin \Anc_G(i)$, and on the other $j \in \Anc_G(k)$, so if $k$ were in $\Anc_G(i)$, it would follow that $j \in \Anc_G(i)$, contradicting $j \notin \Anc_G(i)$.
      Hence, in both cases, we must have $k \not\in \Anc_G(i)$.
      After orienting $i \suh k$ in $H$, the resulting mixed graph still represents $G$.
    \item[$\C{R}3$] ``If $i \suh j \hus k$ in $H$, and $i \suo l \ous k$ in $H$, and $l \suo j$ in $H$, and $i$ and $k$ are not adjacent in $H$, then orient $l \suh j$.''\\
%      \Joris{Note $j \notin J$, and $l \notin J$, and therefore also $i,k\notin J$. So there is no risk of overwriting background knowledge.}
      Since $\lt i, l, k \rt$ is an unshielded triple in $H$ that was not oriented as a collider by $\C{R}0$, we must have that $l \in \Anc_G(\{i,k\})$ by Proposition~\ref{prp:orienting_unshielded_triples}.
      Assume, for the sake of contradiction, that $j \in \Anc_G(l)$. 
      Then $j \in \Anc_G(\{i,k\})$.
      This contradicts that $j \notin \Anc_G(\{i,k\})$ from $i \suh j \hus k$ in $H$.
      Hence, $j \notin \Anc_G(l)$.
      After orienting $l \suh j$ in $H$, the resulting mixed graph still represents $G$.
    \item[$\C{R}4$] ``If $\lt i, q_1, \dots, q_n, j, k \rt$ is a discriminating path in $H$ for $j$, and if $j \ous k$ in $H$, then orient $j \tuh k$ if $j \in \sepset(\{i,k\})$ and orient $q_n \huh j \huh k$ if $j \notin \sepset(\{i,k\})$.''\\ % \sepset(k,i) if ordered!
%      \Joris{Note that $q_1,\dots,q_n \notin J$, $k \notin J$, but $i$ and $j$ might be in $J$. There is the risk of overwriting background knowledge, so we should change the rule (similar to $\C{R}0$). However, since $J\sm\{k,i\}$ is part of $\sepset(k,i)$, there is actually no risk of overwriting background knowledge.}
      It follows immediately from Proposition~\ref{prp:orienting_discriminating_paths} that applying this rule yields an updated mixed graph that still represents $G$.
    \item[$\C{R}8$] ``If $i \tuh j \tuh k$ in $H$, and $i \ouh k$ in $H$, then orient $i \tuh k$.''\\
%      \Joris{$j,k\notin J$, $i$ might be. No risk in overwriting background knowledge.}
      Clearly, $i \in \Anc_G(j)$ and $j \in \Anc_G(k)$ implies $i \in \Anc_G(k)$.
      Thus applying this rule yields an updated mixed graph that still represents $G$.
    \item[$\C{R}9$] ``If $i \ouh k$, and $\pi = \lt i, j, \dots, k \rt$ is an uncovered possibly directed path in $H$ from $i$ to $k$ such that $j$ and $k$ are not adjacent in $H$, then orient $i \tuh k$.''\\
%      \Joris{$j, \dots, k \notin J$, $i \notin J$. No risk to overwrite background knowledge.}
      Suppose $i \notin \Anc_G(k)$.
      The unshielded triple $\lt j, i, k \rt$ was not oriented as a collider by $\C{R}0$, and therefore $i \in \Anc_G(j)$.
      We can repeat such reasoning for each subsequent unshielded triple along the uncovered possibly directed path between $i$ and $k$.
      Overall, by transitivity this implies that $i \in \Anc_G(k)$.
      This, however, contradicts the assumption $i \notin \Anc_G(k)$.
      Therefore, the only possibility is $i \in \Anc_G(k)$.
      Applying the rule therefore yields an updated mixed graph that still represents $G$.
    \item[$\C{R}10$] ``Suppose that $i \ouh k$ in $H$, $j \tuh k \hut l$ in $H$, $\pi_1$ is a uncovered possibly directed path in $H$ from $i$ to $j$, and $\pi_2$ is an uncovered possibly directed path in $H$ from $i$ to $l$. Let $u_1$ be the vertex adjacent to $i$ on $\pi_1$ (possibly $u_1 = j$) and $u_2$ the vertex adjacent to $i$ on $\pi_2$ (possibly $u_2 = l$). If $u_1 \ne u_2$, and $u_1$ and $u_2$ are not adjacent in $H$, then orient $i \tuh k$.''
%      \Joris{$i \notin J$, $k \notin J$, $j \notin J$, $l \notin J$, $u_1 \notin J$, $u_2 \notin J$. No risk to overwrite background knowledge.}
      The unshielded triple $\lt u_1, i, u_2 \rt$ that was not oriented by rule $\C{R}0$ implies that $i \in \Anc_G(u_1)$ or $i \in \Anc_G(u_2)$.
      If $i \in \Anc_G(u_1)$, similar reasoning for each subsequent unshielded triple along the uncovered possibly directed path $\pi_1$ leads us to conclude by transitivity that $i \in \Anc_G(j)$, and since $j \in \Anc_G(k)$, also $i \in \Anc_G(k)$.
      In case $i \in \Anc_G(u_2)$, analogous reasoning for $\pi_2$ leads to the same conclusion, $i \in \Anc_G(k)$.
      In both cases, applying the rule yields an updated mixed graph that still represents $G$.
  \end{description}
  Since each of these orientation rules leaves the skeleton (adjacencies) of $H$ invariant,
  %and each orientation of an edge mark leads to a valid edge,
  and after each orientation rule, $H$ still represents $G$,
  $H$ remains a valid CDPAG throughout the orientation phase
  by Proposition~\ref{prp:dpag_represents_dmg_is_valid}. 
\end{proof}
We have here omitted orientation rules $\C{R}5$--$\C{R}7$ since these are only necessary if selection bias is not ruled out a
priori. 
Those rules can yield edges of the form $\tuo, \out$ and $\tut$ that are not valid in CDPAGs, and require the more general
formalism of PAGs extended with input nodes (which is yet to be written down).


\subsection{Completeness and consistency}

\Joris{TODO: this subsection still needs an overhaul}

\Joris{We should double-check whether in the code we are correctly using the (a)symmetric sepset! Note that $i \isPerp j \given S$ iff $j \isPerp i \given S$  for $i,j \in V$; if $i \in V$ and $j \in J$ then we only consider whether $i \isPerp j \given S$.}

For acyclic $G$, the FCI algorithm was shown to be complete \cite{Zhang2008_AI} in the sense that all edge marks that could possibly be oriented based on the information in $\IM_\sigma(G)$ will be oriented.
Using results on the characterization of Markov equivalence classes of DMAGs \cite{Ali++2005}, it can be additionally shown that the DPAG output by FCI represents the Markov equivalence class of $G$ in case $G$ is acyclic.
These completeness results are rather nontrivial and require very long proofs, and we will therefore not provide these here.
By employing acyclifications, these results have been extended to cyclic $G$ \cite{MooijClaassen_UAI_20}.
We will only state the completeness results here, and refer the interested reader to the original papers for the proofs.

Interpret FCI as a mapping that maps an independence model (on $V$) to a mixed graph (with vertex set $V$).
It maps the independence model of a DMG $G$ to the DPAG $\FCI(\IM_\sigma(G))$.
For two DMGs $G_1,G_2$ with the same vertex set $V$, we will write $G_1 \stackrel{\sigma}{\sim} G_2$ if $\IM_\sigma(G_1) = \IM_\sigma(G_2)$, i.e., if $G_1$ and $G_2$ are $\sigma$-Markov equivalent.
\begin{Thm}\label{thm:fci_sound_complete}
The FCI algorithm is:
  \begin{enumerate}[label=(\roman*)]
%    \item \emph{sound}: for all DMGs $\C{G}$, $\FCI(\IM_\sigma(\C{G}))$ represents $\C{G}$;\label{theo:fci_sound_complete_sound}
    % or (explicit version):} satisfies the following properties:
    % \begin{compactenum}
    %   \item two vertices $i,j$ are adjacent in $\C{P}$ if and only if $i,j$ cannot be $\sigma$-separated by any subset of nodes (not containing $i$ or $j$) in $\C{G}$;
    %   \item if $i \sto j$ in $\C{P}$ (i.e., $i \to j$ in $\C{P}$ or $i \cto j$ in $\C{P}$ or $i \oto j$ in $\C{P}$), then $j \notin \ansub{\C{G}}{i}$;
    %   \item if $i \to j$ in $\C{P}$ then $i \in \ansub{\C{G}}{j}$.
    % \end{compactenum}
    \item \emph{arrowhead complete}: for all DMGs $G$, if $i \notin \Anc_{\tilde{G}}(j)$ for all DMGs $\tilde{G} \stackrel{\sigma}{\sim} G$, then there is an arrowhead $i \hus j$ in $\FCI(\IM_\sigma(G))$;
    \item \emph{tail complete}: for all DMGs $G$, if $i \in \Anc_{\tilde{G}}(j)$ for all DMGs $\tilde{G} \stackrel{\sigma}{\sim} G$, then there is a tail $i \tuh j$ in $\FCI(\IM_\sigma(G))$;
    \item \emph{Markov complete}: for all DMGs $G_1$ and $G_2$, $G_1 \stackrel{\sigma}{\sim} G_2$ if and only if $\FCI(\IM_\sigma(G_1)) = \FCI(\IM_\sigma(G_2))$.
  \end{enumerate}
\end{Thm}
\begin{proof}
  We refer the reader to \cite{MooijClaassen_UAI_20}.
  \begin{comment}

  Sketch:
  The main idea is the following (see also Figure~\ref{fig:different_graphs}). For all DMGs $G$, $\IM_\sigma(G) = \IM_d(G')$ for any acyclification $G'$ of $G$ (Proposition~\ref{prop:Markov_equivalence_acyclification}). 
  Hence FCI maps any acyclification $G'$ of $G$ to the same DPAG $\FCI(\IM_\sigma(G))$, and thereby any conclusion we draw about these acyclifications can be transferred back to a conclusion about $G$ by means of Proposition~\ref{prop:acyclification}. A complete proof is given in Section~\ref{app:proofs} of the Supplementary Material.

  Complete proof:
  %  Note that an alternative formulation of the soundness of FCI in the $\sigma$-separation setting is that the DPAG $\FCI(\IM_\sigma(\C{G}))$ output by FCI represents all DMAGs induced by all acyclifications of $\C{G}$ (see also Figure~\ref{fig:different_graphs}).
  To prove soundness, let $\C{G}$ be a DMG and let $\C{P} = \FCI(\IM_\sigma(\C{G}))$. 
  The acyclic soundness of FCI means that for all ADMGs $\C{G}'$, $\FCI(\IM_d(\C{G}'))$ represents $\C{G}'$. 
  Hence, by Proposition~\ref{prop:Markov_equivalence_acyclification}, $\C{P}$ represents $\C{G}'$ for all acyclifications $\C{G}'$ of $\C{G}$. 
  But then $\C{P}$ must represent $\C{G}$:
  \begin{itemize} 
    \item if two vertices $i,j$ are adjacent in $\C{P}$ then there is a $\sigma$-inducing path between $i,j$ in any acyclification of $\C{G}$, which holds if and only if there is a $\sigma$-inducing path between $i,j$ in $\C{G}$ (Proposition~\ref{prop:acyclification}(\ref{prop:acyclification_$\sigma$-inducing_walk});
    \item if $i \suh j$ in $\C{P}$, then $j \notin \Anc_{G'}(i)$ for any acyclification $\C{G}'$ of $\C{G}$, and hence $j \notin \Anc_{G}(i)$ (Proposition~\ref{prop:acyclification}(\ref{prop:acyclification_anc}));
    \item if $i \tuh j$ in $\C{P}$, then $i \in \Anc_{G'}(j)$ for all acyclifications $\C{G}'$ of $\C{G}$, and hence $i \in \Anc_{G}(j)$ (Proposition~\ref{prop:acyclification}(\ref{prop:acyclification_non_anc})).
  \end{itemize}

  To prove arrowhead completeness, let $\C{G}$ be a DMG and suppose that $i \in \Anc_{G}(j)$ in any DMG $\tilde{\C{G}}$ that is $\sigma$-Markov equivalent to $\C{G}$.
  Since $\C{G}^\acy$ is $\sigma$-Markov equivalent to $\C{G}$, this implies in particular that for all ADMGs $\tilde{\C{G}}$ that are $d$-Markov equivalent to $\C{G}^\acy$, $i \in \Anc_{\tilde{G}}(j)$.
  Because of the acyclic arrowhead completeness of FCI, there must be an arrowhead $i \hus j$ in $\FCI(\IM_d(\C{G}^\acy)) = \FCI(\IM_\sigma(\C{G}))$.
  Tail completeness is proved similarly.

  To prove Markov completeness: Proposition~\ref{prop:Markov_equivalence_acyclification}
  implies both $\IM_\sigma(\C{G}_1) = \IM_d(\C{G}_1^\acy)$ and
  $\IM_\sigma(\C{G}_2) = \IM_d(\C{G}_2^\acy)$. From the acyclic Markov completeness of FCI
  (Proposition~\ref{prop:PAG=MEC} in the Supplementary Material), it then follows that
  $\C{G}_1^\acy$ must be $d$-Markov equivalent to $\C{G}_2^\acy$, and
  hence $\C{G}_1$ must be $\sigma$-Markov equivalent to $\C{G}_2$.

  Alternatively, the statement of this theorem can be seen to be a special case of Theorem~\ref{theo:generalizing_soundness_completeness}, applied with the trivial background knowledge $\Psi(\C{G}) = 1$ for all DMGs $\C{G}$, to $\Phi : \C{G} \mapsto \FCI(\IM_\sigma(\C{G}))$, combined with the known (acyclic) soundness and completeness results of FCI \cite{Zhang2008_AI}. 
  Note here that the trivial background knowledge $\Psi(\C{G}) = 1$ satisfies Assumption~\ref{ass:acybk} as follows immediately from Proposition~\ref{prop:acyclification}.
  \end{comment}
\end{proof}


Arrowhead and tail completeness express that the DPAG output by FCI is maximally oriented: any arrowhead or tail that could possibly be deduced from $\IM_\sigma(G)$, will be oriented.
The soundness and Markov completeness properties together imply that the DPAG $\FCI(\IM_\sigma(G))$ output by FCI, when given as input the $\sigma$-independence model of a DMG $G$, represents the $\sigma$-Markov equivalence class of $G$.
In other words, FCI provides a \emph{graphical characterization} of the $\sigma$-Markov equivalence class of a DMG.

While the soundness of FCI allows us to directly read off some (non-)ancestral relations from the DPAG output by FCI, 
this is not all causal information that is identifiable from the $\sigma$-Markov equivalence class. 
Indirectly, one can also read off additional (non-)ancestral relations, direct causal relations, the absence of confounding, and the absence of cycles. 
For details, see \cite{MooijClaassen_UAI_20}.

In practice, of course we often do not know the independence model $\IM_\sigma(G)$, and have to resort to testing conditional independences in finite samples.
The ``oracle'' soundness and completeness results formulated above, combined with asymptotically consistent conditional independence tests, leads (assuming $i\sigma$-faithfulness) directly to the asymptotic consistency of the FCI algorithm.
We formulate this as the following corollary.
\begin{Cor}[FCI Consistency]
  Let $M$ be a simple SCM with observed variables $O \subseteq V \cup W$, and without exogenous input variables (i.e., $J = \emptyset$).
  Assume that $M$ is $\sigma$-faithful w.r.t.\ $O$.
  Denote the observable causal graph as $G := G^O(M)$.
%  Our goal will be to deduce as much as possible about the observed causal graph $G := G^O(\C{M})$ from the observed distribution $P_{\C{M}}(X_O)$.
  When using asymptotically consistent conditional independence tests on i.i.d.\ samples of the observational distribution $P_{M}(X_O)$, FCI provides an asymptotically consistent estimate $\hat{H}$ of the DPAG $\FCI(\IM_\sigma(G))$ that represents the $\sigma$-Markov equivalence class of $G$. From the estimated DPAG $\hat{H}$, we can read off consistent estimates for:
  \begin{enumerate}[label=(\roman*)]
%  \item the absence of causal cycles according to $\C{M}$ (via Proposition~\ref{prop:noncycles});
%  \item the absence/presence of (possibly indirect) causal relations according to $\C{M}$ (via Propositions~\ref{prop:cdpag_ancestors_dmg} and \ref{prop:cdpag_nonancestors}); 
%  \item the absence of confounders according to $\C{M}$ (via Proposition~\ref{prop:identify_nonconfounders});
%  \item the absence/presence of direct causal relations according to $\C{M}$ (via Proposition~\ref{prop:identifiable_direct_causes}). 
%  \end{compactenum}
  \item the absence/presence of (possibly indirect) causal relations according to $M$; 
  \item the absence of confounding according to $M$;
  \item the absence/presence of direct causal relations according to $M$;
  \item the absence of causal cycles according to $M$.
  \end{enumerate}
\end{Cor}
\begin{proof}
  We refer the reader to \cite{MooijClaassen_UAI_20}.
\end{proof}
Note that this corollary also applies to L-CBNs, as these can be considered to be special cases of simple SCMs.


\Joris{TODO:
\begin{enumerate}
  \item Define FCJ-JCI12r (call it CFCI for conditional FCI; oops, bad idea, that clashes with conservative FCI! Perhaps FCI-WIN instead? "With Input Nodes". Or FCI-C where the C stands for conditional)
  \item Prove its completeness (this might be easy now because we can use Y-structures with dummy variables to mimic the CI). By the way: if we can show that FCI can never orient the edge mark on $j \suh v$ if $j$ is not adjacent to any other nodes, we might get a poor-man's implementation of FCI-CCI by just restricting any edges between the nodes in $J$ in standard FCI.
\end{enumerate}
See file \texttt{old/fci-extensions.tex} in \texttt{jci-paper} repository.
}
